%\Documentstyle[preprint,prb,aps]{revtex}
%\documentstyle[prl,aps]{revtex}
% % TDDFT for non BO systems                  % %
% % Draft 10, 29 March 2001                   % %
% % Last revised by SBT, based on BLW's 9c    % %
%\input{tcilatex}


\documentstyle[amstex,amssymb,prb,12pt,aip,aps]{revtex}
%%%%%%%%%%%%%%%%%%%%%%%%%%%%%%%%%%%%%%%%%%%%%%%%%%%%%%%%%%%%%%%%%%%%%%%%%%%%%%%%%%%%%%%%%%%%%%%%%%%%%%%%%%%%%%%%%%%%%%%%%%%%
%TCIDATA{OutputFilter=LATEX.DLL}
%TCIDATA{LastRevised=Thursday, August 16, 2001 19:23:42}
%TCIDATA{<META NAME="GraphicsSave" CONTENT="32">}
%TCIDATA{Language=American English}
%TCIDATA{CSTFile=revtex.cst}

\include{epsf}
\include{amssymb}
\include{amsfonts}
\epsfverbosetrue 
\newcommand{\be}{\begin{equation}}
\newcommand{\ee}{\end{equation}}
\newcommand{\bea}{\begin{eqnarray}}
\newcommand{\eea}{\end{eqnarray}}
\newcommand{\ba}{\begin{array}}
\newcommand{\ea}{\end{array}}
\newcommand{\half}{1/2}
\newcommand{\pprime}{\prime \prime}
\newcommand{\ppprime}{\prime \prime \prime}
\def\rank{\mathop{\rm rank}\nolimits}
\input{tcilatex}

\begin{document}
\title{Constructive Approach to Time-dependent Density Functional Theory
Beyond the Born-Oppenheimer Approximation I. Contraction Maps and Invariance
Groups for Constrained Search in Dynamics}
\author{B. Weiner \thanks{%
Permanent address Department of Physics, Pennsylvania State University,
Dubois PA 15801, U.S.A.}}
\address{and \\
S.B. Trickey \\
Quantum Theory Project \\
Department of Physics \\
University of Florida \\
Gainesville, FL 32611-8435, U.S.A.\\
May 6, 2001}
\maketitle

\begin{abstract}
\centerline{\bf ABSTRACT} \noindent The usual formulation of time-dependent
density functional theory (TDDFT) relies upon the Runge-Gross argument for
the existence of a 1-1 map between densities and external potentials. Only
in the multi-component plasma generalization due to Li and Tong is this
TDDFT applicable to dynamical nuclei and electrons simultaneously. Both the
original and multi-component versions depend on satisfying certain boundary
conditions and the existence of temporal Taylor expansions, requirements
that have been the subject of skepticism. Little actual application of the
multi-component scheme has occurred.

This paper, the first of a series, extends our alternative formulation of
TDDFT explicitly to the non-Born-Oppenheimer case. The approach does not
depend on the 1-1 map logic as a route to functionals of the density.
Instead, the analysis proceeds by classification of states according to
densities (as in the constrained search approach) into fiber bundles. The
densities then are used as parameters to label states (akin to coherent
state theory), specifically by finding invariances of the fibers and
mappings among them. Formal expressions for functionals of the density that
have a unique value associated to each of these classes of states for all
times are obtained in terms of those transformations. The equation of motion
for the density follows from the time-dependent variational principle.

The subsequent two papers treat (1) approximations to the operators of the
invariance groups to construct approximate functionals systematically and
(2) a corresponding path integral formulation of TDDFT.
\end{abstract}

\newpage

\newpage

%\draft

%\thanks[w] {Permanent address Dept. of Physics, Pennsylvania State
%University, Dubois PA 15801, U.S.A.}  

%\tighten
%\widetext

%\narrowtext

\section{Overview and Motivation}

\noindent Time-dependent density functional theory (TDDFT) is growing in
importance. Diverse examples of its use and development include Refs. %
\citenum{Teil92,Stener95,Saalmann96,Tong97,Serra97a,Ullrich97,AG97,Vignale97,Calvayrac98,Reinhard98,Ullrich98,Telnov98,Tong98,Campillo99,Vasiliev99,Knospe99,Hessler}
as well as papers cited in surveys such as Refs. \citenum{Gross96,Casida95}.
To date less attention has gone to the coupled motions of electrons and
nuclei than to electron excitations and response with static nuclei.

Though a TDDFT for temporally periodic potentials was formulated earlier %
\cite{Bartolotti81,DebGhosh82}, the now-conventional form was introduced by
Runge and Gross \cite{RG84}, hereafter RG. As with the original ground state
Hohenberg-Kohn argument \cite{HK}, RG proceeded by reductio ad absurdum to
the result that, for a $v$-representable time-dependent density, under
certain assumptions the map between the time-dependent external potential
and that density is invertible. For those assumptions, this invertibility
establishes the existence, though not the structure, of a TDDFT. Issues of
boundary conditions on the potentials and densities and implicit assumptions
for which the RG proof is valid arose shortly \cite{Xu85,Dhara87}.

A Levy-Lieb constrained search \cite{Levy79,LevyPerdew85,Lieb} formulation
of the density-potential mapping was given by Kohl and Dreizler \cite{KD86}.
Posed in the context of scattering problems, it did not provide a procedure
for constructing solutions. Then the RG logic was extended to include
time-dependent forces on spatially fixed nuclei \cite{KD86,Raj94A,RajBuot94B}%
. However, coupled electronic and nuclear motion usually is treated via Li
and Tong's extension \cite{LiTong86} to multi-species systems. There is a
considerable variety of other work on the RG logic that, while of general
relevance, is not specific to the issues treated here \cite%
{Qian98,Chermyak00}.

For explicitly wave-function-based theory, a general approach to
non-Born-Oppenheimer dynamics has been given by Deumens, \"Ohrn, and
collaborators \cite{DeumensA,DeumensB}. They use a parameterization via
group coherent states to obtain solutions to the time-dependent
Schr\"odinger equation (TDSE). Here we take up the non-BO problem via
extension of our DFT analogue \cite{BLWSBT99a,BLWSBT00a} of the
Deumens-\"Ohrn scheme. Our formulation involves a generalization to
parameterization of the TDSE by spatially varying functions (in contrast
with the time-dependent complex scalars that they use).

Specifically, the parameterizing functions are the physical densities that
label the density operators via equivalence classes. Density operators are
emphasized in order to exploit the structure of the constrained search
procedure \cite{Levy79,LevyPerdew85}. The most general form of constrained
search must include both pure states and ensembles corresponding to a
specified physical density. For density operators the distinction between
states and ensembles is not a barrier. Furthermore, the search is linear in
the density operators but nonlinear in the states. The density operator
formulation therefore provides a unified way to expose the internal
structure of constrained search in forms exploitable for both analysis and
approximation. See anticipatory remarks along these lines in Refs. %
\citenum{Lieb} (see p.~263) and \citenum{KD86} (just after eq. (5)).

Ref. \citenum{BLWSBT99a} noted that our formulation implicitly is applicable
to the non-BO case. The explicit form is developed in three papers. This one
gives the formal extension as a time-dependent constrained search. First we
generalize Ref. \citenum{BLWSBT99a,BLWSBT00a} from $N$-fermion operators
(and corresponding Hilbert spaces) alone to explicit inclusion of the
nuclear dynamics. Then we develop transformations that constitute the basis
for constructive realizations of constrained search. We give two such
realizations. Both exploit properties of the contraction map from the $N$%
-body density operator to the density. That map confers a density-labeled
fiber bundle structure on the density operators. The kernel of that map
plays a central role in parameterizing density evolution paths. The more
straightforward parameterization is by reference states in each fiber plus
linear combinations of elements of the kernel. The result is an explicit set
of dynamical equations for TDDFT that, however, depends on a constraint
which is difficult to implement. Therefore we develop an alternative
parameterization of the dynamics by constructing Lie groups that map
intra-fiber (i.e. amongst operators with the same density) and Lie groups
that map inter-fiber (one density to another). Application of the
Time-dependent Variational Principle (TDVP) provides equations of motion for
the density in terms of those Lie group parameters that, in the exact case,
is equivalent to the TDSE.

In the second paper, we obtain practical single-particle equations by
restricting the invariance transformations to one- and two-body
contributions in order to produce approximate time-dependent functionals by
systematic refinement. In the limit of $N$-body contributions, the
functionals would be exact. The results are different from the familiar
TDDFT Kohn-Sham equations but play a corresponding role. By a suitable
further approximation the conventional TDDFT equations are recovered in the
static nucleus limit. The third paper uses the invariance transformations to
formulate a path integral version of TDDFT. That provides a convenient means
to treat causality issues, to analyze DFT linear response, and to treat the
relationship of the stationary phase approximation to the TDVP used in the
first two papers.

\section{Runge-Gross TDDFT - Limitations and Counterexamples}

\noindent Some prior results that illuminate our objectives are summarized
here.

First, the Kohl-Dreizler argument \cite{KD86} has limited value for coupled
nuclear and electron motion because of two assumptions the electron
sub-system is assumed to be in a non-degenerate ground state initially.
Non-degeneracy is not a trivial requirement, while the restriction to an
initial ground-state is essential to their use of a minimum principle and
associated infima.

Next, motivation for the Li-Tong multi-component RG TDDFT comes from
considering the original one-component RG TDDFT. For simplicity take $N_n$
nuclei of one species and $N_e$ electrons. RG TDDFT depends on the
assumption that the external dynamic potential 
\begin{equation}
\hat V = \sum_{i}^{N_e} v({\bf r}_{i},t)  \label{VextofT}
\end{equation}
has a temporal Taylor series expansion about the initial time $t=t_{0}$. Let
the initial spatial density be $n({\bf r}, t=t_{0}) \equiv n_{0}({\bf r})$.
Then invertibility of the potential-density map is demonstrated under this
assumption by considering the contrary, to wit, two external potentials $%
v^{\prime}({\bf r},t)$, $v({\bf r},t)$ that yield the same $n$ and differ by
more than a function of time alone. Nonexistence of such duplicate
potentials requires that a certain surface integral vanish. RG argued that
it vanishes for any physically meaningful potential, though they noted that
the Taylor series assumption disqualifies adiabatically switched-on
potentials. van Leeuwen provided an extension of that original argument,
which, however, does not address the issues raised below \cite{Leeuwen99}.

Xu and Rajagopal gave counterexamples in which the surface integral is not
zero \cite{Xu85}. Though Dhara and Ghosh \cite{Dhara87} concluded that those
counterexamples might not be physically realizable, they also pointed out
that the RG proof is not valid for two important categories of systems, the
homogeneous electron gas \cite{HEGexample} and systems with moving bare
nuclei. For the latter, the temporal Taylor series has singularities of
arbitrarily high order \cite{bareTaylor}. The coefficients of the highly
singular terms are time derivatives of $X(t)$ evaluated at $t_{0}$. Those
all vanish for stationary nuclei and the most singular terms in the Taylor
series in that case go as the ordinary nuclear-electron potential. This
problem is ignored in Ref. \citenum{Saalmann96} [at their equation (12)] but
see below.

For the surface integral issue, Gross and Kohn introduced two requirements
on the RG scheme, one physical, the other formal \cite{GrossKohn90}. The
formal constraint concerns the asymptotic behavior, with respect to $r$, of
the Taylor series coefficients of $v({\bf r}_{i},t)$ and is not relevant to
the present discussion. The physical requirement is that ``experimentally
realizable potentials are always produced by {\em real} (this means, in
particular, normalizable) external charge densities $\rho _{ext}({\bf r}t)$
\ldots .'' Clearly this requirement does not address the case of moving bare
nuclei as the external charge densities, even though such delta-function
charges are normalizable in the usual sense. It might be argued that
non-singular distributions were intended. For nuclear motion this means
screening, in practice, a pseudopotential. This is what Refs. %
\citenum{Teil92,Fois88} do in the BO approximation. Refs. %
\citenum{Saalmann96,Knospe99} do the same in practice, hence avoid the error
just noted. Non-singular distributions are fine practically, but do not
solve the basic problem. Restriction to such distributions would force TDDFT
with nuclear motion never to be all-electron. In turn that would lead to the
odd outcome that the time-independent limit of TDDFT would not recover
normal all-electron time-independent DFT (because of the mandatory nuclear
screening).

Thus one arrives at treatment of the nuclei plus electrons via the Li-Tong %
\cite{LiTong86} multi-component plasma version of RG TDDFT. No full DFT
dynamics (i.e., non-Kohn-Sham, even if purely formal) for the Li-Tong scheme
has, to our knowledge, ever been given. Though in principle general, the
nuclear motion in the Li-Tong approach in practice is purely at the KS
level. See for example Gross, Dobson, and Petersilka \cite{Gross96}. Since
KS dynamics is for the {\em model} nuclear system, there is a challenge to
demonstrate correspondence to the dynamics of the {\em physical} system. In
the classical limit, Gross et al.~obtain the correspondence from the
Ehrenfest theorem, with forces on the nuclei from gradients of the TDKS
potentials. Unexplored, as far as we are aware, is the correspondence for
either narrow wavepackets or full-quantum nuclei.\cite{Xchneglect} The
difficulty in analyzing the dynamics lies in the invertibility arguments,
which provide no insight into the structure of the space of time-dependent
densities. Note also that, by construction, the Li-Tong scheme has the same
dependence on convergence of temporal Taylor series expansions as the
original RG scheme, hence cannot include adiabatically switched external
potentials.

Keep in mind that the development given below does not appeal to the
existence of a 1-1 map between external potentials and allowable densities
as a route to functionals of the density. Instead our analysis proceeds by
classification of states according to densities, just as in constrained
search for stationary states. We then use the densities as parameters to
label states in the spirit of coherent state theory. The result is
functionals of the density that have a unique value associated to each of
these classes of states for all times. Equations of motion for the density
follow from the TDVP.\cite{KramSar}

\section{Density Mappings and Fiber Bundles}

\noindent The objective of this section is to show how the densities at any
time can be used to classify density operators into a fiber bundle structure
and to develop exploitable properties of the bundle. The benefit of this
approach shows up in the following sections, where constrained search as a
function of time is converted into well-characterized, explicit
transformations along and among fibers. Systematic approximation of those
transformations then provides a constructive route to approximate
functionals, the topic of paper II.

Without loss of generality we simplify to a single nuclear species (except
when multi-species aspects need to be noted). Let the generic electron and
nuclear coordinates be ${\bf r}$ and ${\bf R}$ respectively and similarly $s$%
, $S$ for the spins. The composite coordinates then will be written as 
\begin{eqnarray}
{\bf x} &\equiv &({\bf r},s)  \nonumber \\
{\bf X} &\equiv &({\bf R}\,,S)  \nonumber \\
{\bf z} &\equiv &({\bf x}\,,{\bf X})
\end{eqnarray}%
The manifolds of $N_{e}$ electronic coordinates (space and spin) and $N_{n}$
nuclei will be denoted as 
\begin{eqnarray}
{\sf x} &=&\{{\bf x}_{1},{\bf x}_{2},\ldots {\bf x}_{N_{e}}\}  \nonumber \\
{\sf X} &=&\{{\bf X}_{1},{\bf X}_{2},\ldots {\bf X}_{N_{n}}\}
\label{manifdef}
\end{eqnarray}

To treat the electronic and nuclear degrees of freedom in a balanced fashion
and, as well, avoid problems inherent in a formulation that decouples
electron space and spin coordinates from the outset (e.g. symmetry
breakings, spin contamination \cite{PerdewSavin,BLWSBT98,SavinSem2}), we
work with the full space-spin electron and nuclear density 
\begin{equation}
\gamma ({\bf x},{\bf X},t)\equiv \gamma ({\bf z},t)  \label{gammadef}
\end{equation}%
rather than the more familiar spatial density $n({\bf r},{\bf \ R},t)$ or
spin-polarized spatial densities. Note that $\gamma $ integrates to the
product $N_{e}N_{n}$, not their sum. $\gamma ({\bf x},{\bf X},t)$ can be
interpreted as the joint probability for the simultaneous occurrence at $t$
of an electron at position ${\bf x}$ and a nucleus at position ${\bf X}$.
BWBW This has the further benefit of avoiding possible problems introduced
by the translational invariance of isolated systems \cite{KreibGross}.BWBW
The single species space-spin densities can be recovered by integration 
\begin{eqnarray}
\gamma _{e}({\bf x},t) &=&\frac{1}{N_{n}}\int \gamma ({\bf x},{\bf X},t)d\,%
{\bf X}  \nonumber \\
\gamma _{n}({\bf X},t) &=&\frac{1}{N_{e}}\int \gamma ({\bf x},{\bf X},t)d\,%
{\bf x}
\end{eqnarray}

With two species, it is easiest to formulate the problem in second
quantization. The field operators are 
\begin{eqnarray}
\Phi^{\dagger}_{e}({\bf x}) & = & \sum_{i} \phi^{*}_{i}({\bf x}%
)a_{i}^{\dagger}  \nonumber \\
\Phi^{\dagger}_{n}({\bf X}) & = & \sum_{j} \phi_{n,j}^{*}({\bf X}%
)b_{j}^{\dagger}  \label{fieldop}
\end{eqnarray}
where the $\phi_{i}({\bf x})$, $\phi_{n,j}({\bf X})$ are elements of
complete one-particle basis sets for electrons and nuclei respectively and $%
a_{i}^{\dagger}$, $b_{j}^{\dagger}$ are electron and nuclear creation
operators respectively, with commutation relations 
\begin{eqnarray}
\left[ a_{i}^{\dagger}, a_{j} \right]_{+} & = & \delta_{ij}  \nonumber \\
\left[ a_{i}^{\dagger}, a^{\dagger}_{j} \right]_{+} & = & \left[ a_{i},
a_{j} \right]_{+} = 0  \nonumber \\
\left[ b_{i}^{\dagger}, b_{j} \right]_{\pm} & = & \delta_{ij}  \nonumber \\
\left[ b_{i}^{\dagger}, b^{\dagger}_{j} \right]_{\pm} & = & \left[
b_{i},b_{j} \right]_{\pm} = 0
\end{eqnarray}
Commutation or anti-commutation depends on the nuclear species (of course)
and all inter-species commutators are zero. Frequently products of these
operators for both an electron and a nucleus occur. A compressed notation
for such products 
\begin{equation}
\Phi^{\dagger}({\bf z})\Phi ({\bf z}) \equiv\Phi^{\dagger}_{e}({\bf x})
\Phi_{e}({\bf x})\Phi^{\dagger}_{n}({\bf X}) \Phi_{n}({\bf X})
\label{fieldopprod}
\end{equation}
will be used.

Generalizing Refs.~\citenum{BLWSBT99a,BLWSBT00a}, the total system density
operator $D^{N}(t)$ is defined in the tensor product Hilbert space of the $%
N_{e}$ electrons and $N_{n}$ nuclei, with $N=N_{e}+N_{n}$. The
time-dependent density then is 
\begin{equation}
\gamma ({\bf z},t)=Tr\left\{ \Phi ^{\dagger }({\bf z})\Phi ({\bf z}%
)D^{N}(t)\right\} =\gamma ({\bf x},{\bf X},t)
\end{equation}%
In the coordinate representation, the total system wave function $\Psi ({\sf %
x},{\sf X},t)$ can be expressed in terms of the vacuum state $|0\rangle $ as 
\begin{eqnarray}
\Psi ({\sf x},{\sf X},t) &=&\langle {\bf x}_{1},\ldots ,{\bf x}_{N_{e}},{\bf %
X}_{1},\ldots ,{\bf X}_{N_{n}}|\Psi (t)\rangle  \nonumber \\
&=&\langle 0|\Phi _{e}({\bf x}_{1})\ldots \Phi _{e}({\bf x}_{N_{e}})\Phi
_{n}({\bf X}_{1})\ldots \Phi _{n}({\bf X}_{N_{n}})\Psi (t)\rangle
\end{eqnarray}%
where $|\Psi (t)\rangle \in {\cal H}^{N_{e}}\otimes {\cal H}^{N_{n}}$ and
the associated $N$-body density operator integral kernels are 
\begin{eqnarray}
&&D^{N}({\bf x}_{1}^{\prime },\ldots {\bf x}_{N_{e}}^{\prime },{\bf X}%
_{1}^{\prime },\ldots {\bf X}_{N_{n}}^{\prime };{\bf x}_{1},\ldots {\bf x}%
_{N_{e}},{\bf X}_{1},\ldots {\bf X}_{N_{n}},t)  \nonumber \\
&=&\Psi ^{\ast }({\bf x}_{1}^{\prime },\ldots {\bf x}_{N_{e}}^{\prime },{\bf %
X}_{1}^{\prime },\ldots {\bf X}_{N_{n}}^{\prime },t)\Psi ({\bf x}_{1},\ldots 
{\bf x}_{N_{e}},{\bf X}_{1},\ldots {\bf X}_{N_{n}},t)
\end{eqnarray}%
(Remark that these objects are called density matrices in the theoretical
chemistry literature, terminology that is inconvenient here.) The
one-electron density operator kernel $D_{e}^{1}({\bf x}^{\prime },{\bf x},t)$
arises from integrating out all the other electron coordinates {\em and all}
the nuclear coordinates. The converse holds for the nuclear density operator
kernel. In the BO approximation, the electron density operator kernel is a
conditional probability relative to the nuclear positions. A relevant
discussion is in Sj\"{o}qvist and Hedstr\"{o}m \cite{Sjoqvist}.

With one substantive alteration, the previous formulation in Refs.~%
\citenum{BLWSBT99a,BLWSBT00a}, involving trace class operators,
Hilbert-Schmidt operators, and contraction maps to one-particle function
spaces goes through essentially unchanged (except for obvious nuclear spin
symmetry details). The alteration arises because it is useful to formulate
the TDVP in terms of unnormalized states and handle the normalization
explicitly.

The central quantity is the contraction map that yields the combined density
from the density operators 
\begin{eqnarray}
\Xi &:&D^{N}(t)\rightarrow \gamma ({\bf z},t)  \nonumber \\
\gamma ({\bf z},t) &=&Tr\{\Phi ^{\dagger }({\bf z})\Phi ({\bf z})D^{N}(t)\}
\label{Xidefn}
\end{eqnarray}%
Here 
\begin{eqnarray}
Tr\{D^{N}(t)\} &=&1\;\;\forall \;t  \nonumber \\
\int \gamma ({\bf z},t)d{\bf z} &=&N_{e}N_{n}\;\forall \;t
\end{eqnarray}%
$\Xi $ itself is a linear map from the space ${\cal B}_{1}({\cal H}^{N})$ of
trace class operators on $N$-particle Hilbert space to absolute integrable
functions on ${\Bbb R}^{6}\times {\Bbb C}^{4}$.

The density labels equivalence classes $[\gamma ]_{N}$ of positive
trace-class operators $Y$ and $W$ by 
\begin{equation}
Y\,,W\in \lbrack \gamma ]_{N}\Leftrightarrow \Xi (Y)=\Xi (W)=\gamma
\end{equation}%
These equivalence classes are elements of a bundle \cite{Bruhat}. By
inspection it is a fiber bundle (specifically an affine bundle) \cite%
{affinebundleref}; see below. Denote it as ${\frak B}_{aff}$. The base of
the bundle is the set of densities, which belong to $L_{1}({\Bbb R}%
^{6}\times {\Bbb C}^{4})$. The elements of a fiber are the members of an
equivalence class which form an affine space; see below. The identification
follows.

If restricted to normalized, positive, trace-class operators, this
classification is physically equivalent to the original constrained search
scheme states are grouped by the density they yield. There are analytical
advantages to working with the density operators instead of the states (the
problem is linear in the density operators but nonlinear in the states). To
gain those advantages it helps to extend consideration, temporarily, to all
trace class operators. Doing so has two consequences that must be addressed
eventually (i) not all the quantities $f$ generated by the contraction map $%
\Xi $ are physical densities $\gamma $ (everywhere positive), (ii) for a
genuine physical density $\gamma $ not every trace-class operator that
contracts to it is positive. With these caveats in mind, the extended
equivalence class is 
\begin{equation}
\lbrack f]_{N}=\left\{ Y\in {\cal B}_{1}({\cal H}^{N});\Xi (Y)=f\right\}
\label{genequiv}
\end{equation}%
(Aside different $f$'s do not necessarily have the same norm.) The advantage
gained is that the kernel of $\Xi $ 
\begin{equation}
\ker \{\Xi \}=\{Y;Tr\,[\Phi ^{\dagger }\Phi Y]=0\}\equiv \left\{ U-W;\Xi
\left( U\right) =\Xi \left( W\right) \right\}  \label{kerdef}
\end{equation}%
is a linear subspace of ${\cal B}_{1}({\cal H}^{N})$. Then every element in
a given equivalence class is expressible as 
\begin{equation}
Y=\sum_{i=1}^{d}\alpha _{i}W_{i}\qquad ;\qquad \sum_{i=1}^{d}\alpha _{i}=1
\label{YexpandW}
\end{equation}%
BWBW which shows that $[f]_{N}$ is an affine space, and BWBW where 
\begin{equation}
W_{i}=Y_{ref}+\xi _{i};\;\;i>0  \label{WifromYref}
\end{equation}%
Here $Y_{ref}$ is a fixed, arbitrary element of the particular $[f]_{N}$
(i.e. is not in $\ker \{\Xi \}$) and the $\{\xi _{i}\}$ constitute a basis
for $\ker \Xi $ \cite{Yref}. Each fiber thus is identified by a reference
element in $[f]_{N}$, while its internal structure is characterized by the
kernel via eqs. (\ref{YexpandW}, \ref{WifromYref}). Eq. (\ref{YexpandW}) can
be recast as 
\begin{eqnarray}
Y &=&Y_{ref}+U  \nonumber \\
U &=&\sum_{i=1}^{d}\alpha _{i}\xi _{i}  \label{YYref}
\end{eqnarray}%
and the bundle is affine \cite{diffbundle}. It follows that the orthogonal
projector onto any pure $N$-particle state $|\Psi _{N}\rangle $ with density 
$\gamma $ can be written as 
\begin{equation}
\frac{\left| \Psi _{N}({\bf c})\right\rangle \left\langle \Psi _{N}({\bf c}%
)\right| }{\left\langle \Psi _{N}({\bf c})|\Psi _{N}({\bf c})\right\rangle }%
=Y_{ref}(\gamma )+\sum_{i=1}^{d}c_{i}\xi _{i}  \label{Psicxi}
\end{equation}%
This expression is key to pure-state constrained search. Changes in values
of the $\{c_{i}\}$ move within a fiber, i.e. search over all system density
operators corresponding to a fixed density. Changing to another fiber means
another density has been selected. Both TDDFT and ground state DFT result by
combination with the appropriate variation principle (stationary action for
TDDFT, Rayleigh-Ritz variation principle for ground state).

\section{Explicit Implementation of Time-dependent Constrained Search}

\noindent The next task is to implement the search along and between fibers
at any time to determine the solution of the time-dependent Schr\"odinger
equation (TDSE). Refs.~\citenum{BLWSBT99a,BLWSBT00a} extended the work of
Kramer and Sarceno \cite{KramSar} (``K\&S'' to distinguish from Kohn and
Sham, commonly denoted ``KS'') to include parameterization of the TDSE via
the density as a function of time and electronic coordinate ${\bf x}$. This
section gives the extension to the density for coupled electron and nuclei
as functions of the coordinate ${\bf z}$ and time $t$ and the resulting
dynamics in terms of the coefficients ${\bf c}$ in eq. (\ref{Psicxi}) and
the density $\gamma$. To exploit the fiber bundle properties the argument
here differs a bit from that of Ref. \citenum{BLWSBT99a}.

A straightforward extension of the original K\&S analysis of the TDSE can be
made for the manifold of states $\left| \Psi (\lambda (t))\right\rangle $
parameterized by real, time-dependent functions $\lambda (t,{\bf z})$. For
these states, the solutions of the TDSE, denoted $\left| \Psi (\lambda
_{\#}(t))\right\rangle $, are characterized by parameterizing functions that
satisfy 
\begin{equation}
\int \eta (t,\,{\bf y},\,{\bf y}^{\prime })\frac{d\lambda ({\bf y}^{\prime })%
}{dt}d{\bf y}^{\prime }=\frac{\delta {\cal E}(t)}{\delta \lambda ({\bf y})}%
\qquad \qquad \forall \;{\bf y}  \label{dEdlambda}
\end{equation}%
where 
\begin{equation}
\eta (t,\,{\bf y},\,{\bf y}^{\prime })=i\left\{ \frac{\delta }{\delta \,%
\tilde{\lambda}({\bf y})}\frac{\delta }{\delta \,\lambda ({\bf y}^{\prime })}%
\,-\,\frac{\delta }{\delta \,\lambda ({\bf y})}\frac{\delta }{\delta \,%
\tilde{\lambda}({\bf y}^{\prime })}\right\} \ln \left\langle \Psi (\tilde{%
\lambda}(t))|\Psi (\lambda (t))\right\rangle \,|_{\tilde{\lambda}=\lambda }
\label{eta}
\end{equation}%
and 
\begin{equation}
{\cal E}(t)=\frac{\left\langle \Psi (\lambda (t))|H|\Psi (\lambda
(t))\right\rangle }{\left\langle \Psi (\lambda (t))|\Psi (\lambda
(t))\right\rangle }  \label{Eoft}
\end{equation}%
Note the explicit norm. The metric (in the space of parameterizing
functions) satisfies 
\begin{equation}
\eta (t,\,{\bf y},\,{\bf y}^{\prime })=-\eta (t,\,{\bf y}^{\prime },\,{\bf y}%
)
\end{equation}

Eq. (\ref{dEdlambda}) can be cast in a form analogous with classical
mechanics, to wit 
\begin{eqnarray}
\frac{d\lambda ({\bf y}^{\prime })}{dt} &=&\int \,\eta (t,\,{\bf y}^{\prime
},\,{\bf y})^{-1}\frac{\delta {\cal E}(t)}{\delta \lambda ({\bf y})}\,d{\bf y%
}\quad \forall \;{\bf y}^{\prime }  \nonumber \\
\dot{\lambda}({\bf y}^{\prime }) &=&\{\lambda ({\bf y}),{\cal E}\}
\label{genEOM}
\end{eqnarray}%
with the generalized Poisson bracket 
\begin{equation}
\{f(t),g(t)\}\equiv \int \frac{\delta f(t)}{\delta \,\lambda ({\bf y})}\zeta
(t,\,{\bf y},\,{\bf y}^{\prime })\frac{\delta g(t)}{\delta \,\lambda ({\bf y}%
^{\prime })}d{\bf y}\,d{\bf y}^{\prime }  \label{genPoisson}
\end{equation}%
and the definition 
\begin{equation}
\zeta (t,\,{\bf y},\,{\bf y}^{\prime })\equiv \eta ^{-1}(t,\,{\bf y},\,{\bf y%
}^{\prime })  \label{etainverse}
\end{equation}%
all under the assumption that $\eta ^{-1}$ exists. This is not a strongly
restrictive assumption \cite{etainv}.

{\em If} the manifold of states $\left| \Psi (\lambda (t))\right\rangle $ is
all of ${\cal H}^{N}$, {\em then} solutions of Eqs. (\ref{genEOM}), (\ref%
{genPoisson}) are equivalent to solutions of the full TDSE except for an
overall phase that also can be determined within the parameterization \cite%
{KramSar,BLW94}. It follows from eq. (\ref{Psicxi}) that, indeed 
\begin{equation}
{\cal H}^{N}=\left\{ \left| \Psi _{N}({\bf c},\gamma )\right\rangle \right\}
\end{equation}%
with an obvious change in notation to account for the $\gamma $ dependence.
Note that the set of coefficients ${\bf c}$ must be constrained such that
the result in eq. (\ref{Psicxi}) is a pure-state projector, in symbols 
\begin{equation}
\{c_{i}\}\subset \Lambda _{Y_{ref}(\gamma )}\ni Y={\rm pure \, state\,
density \,operator}  \label{posconstraint}
\end{equation}
The problem of characterizing the set $\Lambda $ (or finding an alternative)
is deferred to the next section. Proceeding, the equation of motion (\ref%
{genEOM}) that is fully equivalent to the TDSE becomes 
\begin{equation}
\{({\bf c,}\gamma ),{\cal E}\}=(\dot{{\bf c,}}\dot{\gamma})
\label{cgammaeom}
\end{equation}%
Since all the quantities involved are functionals of $\gamma $, this
equation demonstrates the existence of a TDDFT independent of the RG logic.
The claim becomes evident upon partitioning eq. (\ref{cgammaeom}) between
the ${\bf c}(t)$ and $\gamma (t)$. Details are in Appendix A. The outcome is
that the density dynamics can be displayed explicitly, at least formally,
with the $c_{i}$ as functionals of $\gamma $. In terms of them, the exact
temporal path for the density is the solution of 
\begin{equation}
\dot{\gamma}(t)=\tilde{{\cal N}_{1}}(\gamma (t))\left[ \nabla _{\gamma }\,%
{\cal E}(\gamma (t))-\tilde{{\cal N}_{2}}(\gamma (t))\right]
\label{gammadot1}
\end{equation}%
where $\tilde{{\cal N}_{1}}(\gamma (t))$, $\tilde{{\cal N}_{2}}(\gamma (t))$
are density functionals determined by the solutions of the equation of
motion for the coefficients ${\bf c}$, hence self-consistently. See Appendix
A.

\section{Dynamics Parameterized by Group Properties of the Fiber Bundle}

Parameterization via the ${\bf c}$ coefficients has the merit of a clear,
straightforward linkage to the structure of constrained search. There seems
to be no straightforward way, however, to implement the constraint eq. (\ref%
{posconstraint}) requiring the $c_{i}$ to preserve positivity of the
operator $Y$ in eq. (\ref{YYref}). Observe that this constraint is in
addition to having to select an approximation to model the states $\left|
\Psi _{N}({\bf c},\gamma )\right\rangle $. A suggestive aspect of the ${\bf c%
}$ parameterization is that it is not complex analytic in general. A complex
analytic parameterization is technically useful for dealing with the
positivity constraint, as shown in this Section. Such a parameterization
with independent real and imaginary parts always is equivalent to a real
parameterization, hence does not violate the assumptions that give eqs. (\ref%
{dEdlambda}) - (\ref{Eoft}).

Thus we develop an alternative parameterization based on group properties of
transformations of a vector bundle, ${\frak B}_{vect}$, that is isomorphic
to the affine bundle, ${\frak B}_{aff}$, constructed in Section 3. This does
lead to a complex analytic description. Use of this vector bundle also makes
the group theoretic analysis less demanding because we only need to consider
transformations of vector spaces rather than affine spaces. The vector
bundle ${\frak B}_{vect}$ is based on the direct sum decomposition of ${\cal %
B}_{1}\left( {\cal H}^{N}\right) $ given by 
\begin{equation}
{\cal B}_{1}\left( {\cal H}^{N}\right) =\ker \Xi \oplus \ker \Xi ^{c}
\label{directsum}
\end{equation}%
It has a base space $\ker \Xi ^{c}$, the vector space complement of $\ker
\Xi $ in ${\cal B}_{1}\left( {\cal H}^{N}\right) $, which is isomorphic to $%
L_{1}({\Bbb R}^{6}\times {\Bbb C}^{4})$, and fibers $\ker \Xi $ that are all
identical. Note therefore that the base of ${\frak B}_{vect}$ is isomorphic
with the set of absolutely integrable functions that contain the set of
densities and that each equivalence class $[f]_{N}$ is comprised of the
fiber plus a particular base element. Both the fibers and the base are
vector spaces.

Transformations of ${\cal B}_{1}\left( {\cal H}^{N}\right) $ that are
block-triangular with respect to bases adapted to this decomposition, that
is, are of the form 
\begin{equation}
\left( 
\begin{array}{cc}
\left| \ker \Xi \right) \left( \ker \Xi \right| & \left| \ker \Xi \right)
\left( \ker \Xi ^{c}\right| \\ 
0 & \left| \ker \Xi ^{c}\right) \left( \ker \Xi ^{c}\right|%
\end{array}%
\right)  \label{blockxform}
\end{equation}%
induce vector bundle maps ${\frak B}_{vect}\rightarrow {\frak B}_{vect}$ and
form the group, ${\frak G}_{\Xi }$, of automorphisms of ${\frak B}_{vect}$ %
\cite{Bruhat}. Here $\left| A\right) \left( B\right| $ denotes mappings from
the subspace $B$ to the subspace $A$ and the lower block of eq. (\ref%
{blockxform}) is zero because mappings from the fiber to the base destroy
the vector bundle properties.

Usually invariance groups characterizing a property of a class of states are
constructed from unitary transformations. In the present case there is no
unitary group of transformations that operates homogeneously, (i.e. produces
all pure states from any given reference pure state), both in ${\cal H}^{N}$
and the set of pure states contained in each equivalence class $\left[
\gamma \right] _{N}$ and also preserves the vector bundle structure of $%
{\frak B}_{vect}$. However one can construct a commutative, but non unitary,
Lie group, ${\frak G}_{\ker }$, of transformations contained in the general
linear group, $GL\left( {\cal H}^{N}\right) $, of ${\cal H}^{N}$ that does
have these properties. This group induces positive congruence maps of ${\cal %
B}_{1}\left( {\cal H}^{N}\right) \mapsto {\cal B}_{1}\left( {\cal H}%
^{N}\right) $, hence maps states to states (see Appendix B for details). The
group ${\frak G}_{\ker }$ contains subgroups ${\frak G}_{den}$ and ${\frak G}%
_{fib}$ that can be used to give parameterizations of $\ker \Xi ^{c}$ and $%
\ker \Xi $ respectively.

The parameterization of the base $\ker \Xi ^{c}$ that we produce is non
redundant (i.e. one set of parameters corresponds to one element and one
element only) and is given in terms of the group ${\frak G}_{den}\subset 
{\frak G}_{\ker }$ which is represented by transformations of the form 
\begin{equation}
\left( 
\begin{array}{cc}
t\left( T_{B}\right) & 0 \\ 
0 & T_{B}%
\end{array}%
\right)
\end{equation}%
acting on a reference element in $\ker \Xi ^{c}$, where the top block is
determined by the bottom one. In the case of the fibers, $\ker \Xi $, we
produce a non redundant parameterization of {\it pure states }given in terms
of the group ${\frak G}_{fib}$ of transformations 
\begin{equation}
\left( 
\begin{array}{cc}
T_{F} & 0 \\ 
0 & I_{B}%
\end{array}%
\right)
\end{equation}%
acting on a pure state reference, where the bottom block is the identity
operation on the base. A more detailed discussion of the properties of these
transformations, which are determined by integral kernels of operators
acting in ${\cal H}^{N}$, that are diagonal in the coordinate
representation, are given later in this section and in the Appendices.

States (pure or ensemble) of a physical system correspond to positive trace
class operators. In the case of pure states, these operators correspond to
elements of $\ {\cal H}^{N}$. As we are interested in mappings between
states, positive maps of such operators are critical. Since all physical
densities can be associated with pure states, we focus on them. Relevant
details of positive operators and maps are summarized in Appendix C. In
particular, the only two possibilities for mappings between pure states are
unitary transformations and the more general congruence transformations\cite%
{congruence} of ${\cal B}_{1}\left( {\cal H}^{N}\right) $. To incorporate
both types, we therefore consider the most general possible invertible
transformations on ${\cal H}^{N}$, namely those of the General Linear group $%
GL({\cal H}^{N})$ 
\begin{equation}
T\,\in \,GL({\cal H}^{N})\ni T|\Psi \rangle =|\chi \rangle \text{ and }%
T^{-1}|\chi \rangle =|\Psi \rangle \quad \forall |\Psi \rangle ,|\chi
\rangle \in {\cal H}^{N}
\end{equation}%
Because norms appear explicitly in the TDVP, the non-unitarity of these maps
is not a difficulty. It follows that 
\begin{equation}
T=e^{A}\,e^{iB}  \label{eAeiB}
\end{equation}%
with $A^{\dagger }=A$ and $B^{\dagger }=B$ (see Appendix C). All pure state
density operators are mapped to one another by such transformations; see
Appendix D on this point as well as for the basis of a systematic extension
of the present analysis to ensemble density operators.

The most general group, ${\frak G}_{\ker }\subset {\frak G}_{\Xi }\subset
GL\left( {\cal B}_{1}\left( {\cal H}^{N}\right) \right) $, of positive
transformations of ${\cal B}_{1}\left( {\cal H}^{N}\right) $ that both
preserve the fiber bundle structure of ${\frak B}_{vect}$ and also are
congruence transformations then has the following property 
\begin{equation}
{\frak G}_{\ker }=\left\{ T\in GL({\cal H}^{N});T{\sl r}\left\{ \Phi
^{\dagger }\Phi TYT^{\dagger }\right\} =0\;\forall Y\,\in \ker \Xi \text{ }%
\right\}  \label{GLkerXi}
\end{equation}%
By resort to the coordinate eigenfunction representation, we show in
Appendices E and F that only certain subgroups of the commutative subgroup $%
{\cal C}({\cal H}^{N})$ of diagonal coordinate multi-particle
transformations $T(\Upsilon )$ have this property. In terms of the field
operators $\left\{ \Phi ^{\dagger }\left( \kappa \right) ,\Phi \left( \kappa
\right) ;\kappa =\left( {\bf x},{\bf X}\right) \right\} $, these diagonal
transformations are 
\begin{eqnarray}
{\cal C}({\cal H}^{N}) &=&\left\{ T(\Upsilon );T(\Upsilon )YT^{\dagger
}(\Upsilon )\equiv \widehat{T}(\Upsilon )(Y)\right. \,;  \nonumber \\
T(\Upsilon ) &=&\left. e^{\int \Upsilon (\kappa _{1},\ldots ,\kappa
_{N})\Phi ^{\dagger }(\kappa _{1})\Phi (\kappa _{1})\ldots \Phi ^{\dagger
}(\kappa _{N})\Phi (\kappa _{N})}d\kappa _{1}\ldots d\kappa _{N}\right\} 
\nonumber \\
&&\Upsilon (\kappa _{1},\ldots ,\kappa _{N})\in {\Bbb C}
\end{eqnarray}%
Here and in the Appendices the coordinate strings are in the compressed
notation 
\begin{equation}
\kappa _{1}\ldots \kappa _{N_{e}},\kappa _{N_{e}+1}\ldots \kappa _{N}\equiv 
{\bf x}_{1}\ldots {\bf x}_{N_{e}},{\bf X}_{1}\ldots {\bf X}_{N_{n}}
\label{compressed}
\end{equation}%
By construction ${\cal C}({\cal H}^{N})$ has the property 
\begin{equation}
\left[ T(\Upsilon ),\Phi ^{\dagger }\Phi \right] =0
\end{equation}%
Note that {\em all} states (unnormalized in general) belonging to ${\cal H}%
^{N}$ can be expressed as 
\begin{equation}
|\Psi (\Upsilon )\rangle =T(\Upsilon )|\Psi \rangle
\end{equation}%
for any fixed reference state $|\Psi \rangle \in {\cal H}^{N}$; see Appendix
H. Separation of $\Upsilon $ into real and imaginary parts as in eq. (\ref%
{eAeiB}) gives 
\begin{eqnarray}
T(\Upsilon ) &=&e^{\int {\cal V}(\kappa _{1},\ldots ,\kappa _{N})\Phi
^{\dagger }(\kappa _{1})\Phi (\kappa _{1})\ldots \Phi ^{\dagger }(\kappa
_{N})\Phi (\kappa _{N})d\kappa _{1}\ldots d\kappa _{N}}  \nonumber \\
&&\times e^{i\int \omega (\kappa _{1},\ldots ,\kappa _{N})\Phi ^{\dagger
}(\kappa _{1})\Phi (\kappa _{1})\ldots \Phi ^{\dagger }(\kappa _{N})\Phi
(\kappa _{N})d\kappa _{1}\ldots d\kappa _{N}}  \nonumber \\
&\equiv &e^{\breve{{\cal V}}}e^{i\breve{\omega}}  \label{TVomega}
\end{eqnarray}

The group of transformations, ${\cal G}_{\Xi }$, generated by the imaginary
part of $T$ is a unitary subgroup of ${\cal C}({\cal H}^{N})$ that leaves $%
\ker \Xi $ invariant 
\begin{eqnarray}
\!{\cal G}_{\Xi } &=&\left\{ U(\omega );\hat{U}(\omega )(Y)\equiv U(\omega
)YU^{\dagger }(\omega );\right. \omega (\kappa _{1},\ldots ,\kappa _{N})\in 
{\Bbb R}  \nonumber \\
U(\omega ) &=&\left. e^{i\,\int \omega (\kappa _{1},\ldots ,\kappa _{N})\Phi
^{\dagger }(\kappa _{1})\Phi (\kappa _{1})\ldots \Phi ^{\dagger }(\kappa
_{N})\Phi (\kappa _{N})d\kappa _{1}\ldots d\kappa _{N}}\right\}
\label{GXidefn1}
\end{eqnarray}%
Elements of this subgroup map each fiber to itself i.e. are intra fiber maps 
\begin{equation}
Tr\{\Phi ^{\dagger }\Phi U(\omega )YU^{\dagger }(\omega )\}=Tr\{\Phi
^{\dagger }\Phi Y\}\quad \forall \,Y\in {\cal B}_{1}({\cal H}^{N})
\end{equation}%
In fact ${\cal G}_{\Xi }$ is a subgroup of {\em the} invariance group of the
contraction map $\Xi $, since 
\begin{equation}
\Xi \left( U(\omega )YU^{\dagger }(\omega )\right) =\Xi \left( Y\right) 
\text{ }\forall Y\in {\cal B}_{1}\left( {\cal H}^{N}\right)
\end{equation}%
The group ${\frak G}_{\ker }$ can be expressed as the direct product ${\cal G%
}_{\ker }\times {\cal G}_{\Xi }$, where the subgroup, ${\cal G}_{\ker }$, is
generated by ${\cal V}$ (the real part of $\Upsilon $).

To construct a representation of this subgroup, which contains both inter
and intra fiber maps, restrictions must be imposed upon ${\cal V}(\kappa
_{1},\ldots ,,\kappa _{N})$ so that eq. $\left( \ref{GLkerXi}\right) $ is
satisfied. In Appendix F we show that the diagonal integral kernel $T_{d}$
of this transformation may be written as 
\begin{equation}
|T_{d}(\kappa _{1},\ldots \kappa _{N})|^{2}=1+\phi (\kappa _{1},\ldots
\kappa _{N})=e^{2{\cal V}(\kappa _{1},\ldots ,,\kappa _{N})}
\end{equation}%
(note Appendix F, eq. (\ref{phidef})), i.e., $T_{d}=e^{\breve{{\cal V}}%
}=\left( 1+\breve{\phi}\right) ^{\frac{1}{2}}$. This decomposition is
significant because of the requirement (see Appendix F) that $\breve{\phi}$
needs to be strongly orthogonal to $K_{D}$, the diagonal part of $\ker \Xi $
in the coordinate representation. In other words, the integral kernel $\phi
(\kappa _{1},\ldots \kappa _{N})$ of $\breve{\phi}$ must be strongly
orthogonal to the integral kernels of all diagonal operators that lie in $%
\ker \Xi $. This condition is obtained, (see Appendix F), by considering the
decomposition of $Y\in {\cal B}_{1}({\cal H}^{N})$ as 
\begin{equation}
Y=Y_{K_{OD}}+Y_{K_{D}}+\Gamma \left( \gamma \right)
\end{equation}%
where we decompose $\ker \Xi $ into off diagonal and diagonal parts, i.e. $%
\ker \Xi $ $=K_{OD}\oplus K_{D}$, and $Y_{K_{OD}}\in K_{OD}$, $Y_{K_{D}}\in
K_{D}$, and $\Gamma \left( \gamma \right) \in \ker \Xi ^{c}$. Observe that
the transformed diagonal part $\left( 1+\breve{\phi}\right) ^{\frac{1}{2}%
}Y_{K_{D}}\left( 1+\breve{\phi}\right) ^{\frac{1}{2}}=\left( 1+\breve{\phi}%
\right) Y_{K_{D}}$ only belongs to $Y_{K_{D}}$ if $\phi \left( \kappa
_{1},\ldots \kappa _{N}\right) $ is strongly orthogonal to $Y_{K_{D}}\left(
\kappa _{1},\ldots \kappa _{N}\right) $, while $\left( 1+\breve{\phi}\right)
^{\frac{1}{2}}Y_{K_{OD}}\left( 1+\breve{\phi}\right) ^{\frac{1}{2}\dagger }$
always belongs to $K_{OD}$ (note that $\left( 1+\breve{\phi}\right) ^{\frac{1%
}{2}}$ does not commute with $Y_{K_{OD}}$) and $\left( 1+\breve{\phi}\right)
\Gamma \left( \gamma \right) $ belongs to $\ker \Xi ^{c\text{ }}$ both in
the case when $\breve{\phi}$ is strongly orthogonal to $\Gamma \left( \gamma
\right) $ and when it is not. Clearly this map maintains the vector bundle
structure of ${\frak B}_{vect}$ as it does not send elements of the fibers
to the base $\ker \Xi ^{c}$.

Not all elements of ${\cal G}_{ker}$ are intra-fiber maps, as in general $%
\left( 1+\breve{\phi}\right) \Gamma \left( \gamma \right) =\Gamma \left(
\gamma ^{\prime }\right) +Y_{K_{D}}^{\prime }$, and thus $\left( 1+\breve{%
\phi}\right) $ is an intra-fiber map only if $\gamma ^{\prime }=\gamma $. In
terms of eq. (\ref{YYref}), the issue is that elements of ${\cal G}_{ker}$
may change $Y_{ref}\rightarrow Y_{ref}^{\prime }$, where $Y_{ref}^{\prime }$
is in a different fiber to $Y_{ref}$. However, ${\cal G}_{ker}$ does contain
a subgroup, ${\cal G}_{fib\text{,}}$ of intra fiber maps. They are treated
in Appendix F also. An essential result is that the ${\cal \breve{V}}_{fib}$
that generate this group are expressed in terms of $\breve{\phi}$ that are
strongly orthogonal to $K_{D}$ and act as the zero operator on $S_{D}$.
Together with the group ${\cal G}_{\Xi }$ the group ${\cal G}_{fib}$ forms,
via direct product, the group ${\frak G}_{fib}$ of intra fiber super
multiplication operators, acting in ${\cal B}_{1}\left( {\cal H}^{N}\right) $
mapping $\ker \Xi \mapsto \ker \Xi ,$ generated by multiplication operators
(in the coordinate representation) acting in ${\cal H}^{N}$ i.e. ${\frak G}%
_{fib}={\cal G}_{fib}\times $ ${\cal G}_{\Xi }$. This group can be used to
give a non redundant parameterization of $\ker \Xi $, as we show in Appendix
G.

A representation of the group, ${\frak G}_{den}$, of maps between fibers,
that is contained in ${\cal G}_{ker}$, can be constructed (see Appendix F)
using ${\cal \breve{V}}_{den}$ based on $\breve{\phi}$ that are strongly
orthogonal to all diagonal operators that lie in $\ker \Xi $ but are not
strongly orthogonal to $S_{D}$. A non redundant parameterization of $S_{D}$
based on a specific (but arbitrary) reference state $D_{ref}^{N}\left(
\gamma _{0}\right) $ can be formed using this representation (see Appendix
G). It is possible to choose the reference state $D_{ref}^{N}\left( \gamma
_{0}\right) $ and all the states generated from it, (one for each fiber), to
be Independent Particle States or to be states of any prescribed form.

As a result, a non redundant parameterization of the set of pure states in $%
{\frak B}_{vec}$ can be given in terms of the action of the product ${\frak G%
}_{fib}\times {\frak G}_{den}={\cal G}_{fib}\times {\cal G}_{\Xi }\times 
{\frak G}_{den}$ $\equiv {\frak G}_{pure}$ on a pure state reference
element. Elements of ${\cal G}_{fib}\times {\cal G}_{\Xi }$ are intra-fiber
maps, while elements of ${\frak G}_{den}$ are inter-fiber maps. The state $%
D_{ref}^{N}\left( {\cal V}_{den}\right) \equiv D_{ref}^{N}\left( {\cal %
\gamma }\right) $ obtained through the non redundant action of ${\frak G}%
_{den}$ on $D_{ref}^{N}\left( \gamma _{0}\right) $ can be used as the
reference state for the non redundant parameterization of the pure states
contained in $\left[ \gamma \right] _{N}$ through the action of ${\frak G}%
_{fib}$ on $D_{ref}^{N}\left( {\cal \gamma }\right) $, that gives $%
D^{N}\left( {\cal V}_{fib}\left( \gamma \right) ,{\cal \omega }\right) $.
Note that these maps are rank-preserving. The physical significance is that
while these maps do {\em not} explore all elements of every fiber, they {\em %
do} explore all pure-state elements of every fiber if the initial reference
is a pure state. Technical details are in Appendices F and G. In the
following diagram, we collect the primary relationships described in this
section. The top line is the domain, the bottom line the range of the
mappings shown in the middle.

\begin{gather}
\underbrace{{\frak B}_{aff}\backsimeq {\frak B}_{vect}\backsimeq {\cal B}%
_{1}\left( {\cal H}^{N}\right) \equiv \ker \Xi \oplus \ker \Xi ^{c}}%
\mathstrut \medskip  \nonumber \\
\left\downarrow 
\begin{array}{c}
{\frak G}_{\Xi }\subset GL\left( {\frak B}_{1}\left( {\cal H}^{N}\right)
\right) \\ 
\cup \\ 
{\frak G}_{\ker }={\cal G}_{\ker }\times {\cal G}_{\Xi } \\ 
\cup \\ 
{\frak G}_{pure}={\frak G}_{fib}\times {\frak G}_{den} \\ 
{\LARGE \parallel } \\ 
{\cal G}_{fib}\times {\cal G}_{\Xi }\times {\frak G}_{den}%
\end{array}%
\right\downarrow  \nonumber \\
\overbrace{{\frak B}_{aff}\backsimeq {\frak B}_{vect}\backsimeq {\cal B}%
_{1}\left( {\cal H}^{N}\right) \equiv \ker \Xi \oplus \ker \Xi ^{c}} 
\nonumber \\
{\frak G}_{\ker }\subset {\cal C}({\cal H}^{N})\cap \text{Congruence
Transforms of }{\cal B}_{1}\left( {\cal H}^{N}\right)
\end{gather}

\section{Equations of Motion}

Now we use the parameterization of pure states in ${\cal H}^{N}$ provided by
the factorization of ${\frak G}_{pure}$ in terms of ${\cal G}_{\Xi }$, $%
{\cal G}_{fib}$, and ${\cal G}_{den}$ to express the TDVP in terms of
functionals of the physical density in a form amenable to sequential
approximation. First a reference state $|\Psi _{0}\rangle \in {\cal H}^{N}$
must be selected. Such a state could be of a simple form (e.g. an
Independent Particle State) that leads to straightforward expressions for
expectation values. Or, it could be a correlated state (e.g., the product of
an electronic Antisymmetrized Geminal Power State \cite{Ortiz81} and a
variable width wave packet for nuclear positions) that provides a better
basis for approximations - though at the expense of more complicated
expressions for expectation values. Given such a reference state $|\Psi
_{0}\rangle $ and its associated density $\gamma _{0}$, {\em any} pure state
(in general, unnormalized) in ${\cal H}^{N}$ can be parameterized in terms
of the foregoing group factorization as 
\begin{equation}
|\Psi ({\cal V}_{fib}(\gamma ),\omega ,{\cal V}_{den}\rangle =e^{_{{\cal 
\breve{V}}_{fib}(\gamma )}}\,e^{i\breve{\omega}}e^{\breve{{\cal V}}%
_{den}}|\Psi _{0}\rangle  \label{vv1omega1}
\end{equation}%
where 
\begin{eqnarray}
e^{_{{\cal \breve{V}}_{fib}(\gamma )}} &\in &{\cal G}_{fib}  \nonumber \\
e^{\breve{{\cal V}}_{den}} &\in &{\frak G}_{den}  \nonumber \\
e^{i\breve{\omega}} &\in &{\cal G}_{\Xi }  \label{vv1omega2}
\end{eqnarray}%
$\gamma _{0}$ is the density of the state $|\Psi _{0}\rangle $, $e^{\breve{%
{\cal V}}_{den}}$ produces states of different density, $e^{_{{\cal \breve{V}%
}_{fib}(\gamma )}}$ produces states with the same density as $e^{\breve{%
{\cal V}}_{den}}|\Psi _{0}\rangle $ and the shorthand notation for
exponentiated integral operators was introduced at eq.~(\ref{TVomega}).

Equations (\ref{vv1omega1}), (\ref{vv1omega2}) make clear the distinction
between the time propagation of the density $\gamma (t)$ and the states. In
particular, it is vital that different $\breve{{\cal V}}_{den}$'s produce
states with different densities and that all possible densities $\gamma $
associated with pure states can be generated by varying ${\cal V}_{d}(\gamma
_{0})$ such that there is a 1-1 correspondence between $\breve{{\cal V}}%
_{den}$ and $\gamma $. That is, for each $\gamma $ there is a unique
transformation $e^{\breve{{\cal V}}_{den}}\in {\frak G}_{den}$ that maps $%
|\Psi _{0}\rangle \mapsto |\Psi _{\gamma }\rangle $, where $|\Psi _{\gamma
}\rangle $ is {\em some} state that yields $\gamma $. Details are in
Appendices D and G. Parameterization by $\breve{{\cal V}}_{den}$ therefore
is equivalent to parameterization with respect to $\gamma $. Equivalently 
\begin{equation}
|\Psi ({\cal V}_{fib}(\gamma ),\omega ,\gamma )\rangle \equiv |\Psi ({\cal V}%
_{fib}(\gamma ),\omega ,{\cal V}_{den}\rangle
\end{equation}%
Notice in this last expression that the roles of ${\cal V}_{fib}(\gamma )$
and $\omega $ are to determine the time propagation of the states (hence
also the density operators) associated with $\gamma $. The ``fiber factor'' 
\begin{equation}
e^{_{{\cal \breve{V}}_{fib}(\gamma )}}\,e^{i\breve{\omega}}\equiv e^{\breve{%
\Upsilon}}  \label{fiberfactor}
\end{equation}%
selects among elements of the fiber, while $e^{\breve{{\cal V}}_{den}}$ is
the ``base factor'' that determines which $\gamma $ will be reached from the
initial one. From this perspective, the task of constructing density
functionals is equivalent to the task of finding the optimal ${\cal V}%
_{fib}(\gamma )$ and $\omega $ for the particular $\gamma $. By comparison
with the development leading to eq. (\ref{genEOM}), the equation of motion
in terms of this parameterization is 
\begin{equation}
\{({\cal V}_{fib}(\gamma ),\omega ,\gamma ),{\cal E}\}=({\cal \dot{V}}%
_{fib}(\gamma ),\dot{\omega},\dot{\gamma})  \label{Vomegagammaeom}
\end{equation}

In Appendix I we show how to utilize the $\left( {\cal V}_{fib}(\gamma
),\omega \right) $ parameterization in a form which emphasizes the complex
analytic character of the fiber parameterization via use of the complex
fiber parameter 
\begin{equation}
\Upsilon ={\cal V}_{fib}(\gamma )+i\omega  \label{Ffactor}
\end{equation}%
The result of this parameterization is a different EOM, again a pure density
functional dynamical expression, of the same form as eq. (\ref{gammadot1}) 
\begin{equation}
\dot{\gamma}(t)={\cal N}_{1}(\gamma (t))\left[ \nabla _{\gamma }\,{\cal E}%
(\gamma (t))-{\cal N}_{2}(\gamma (t))\right]  \label{gammadot2}
\end{equation}%
but with different functionals and with the positivity constraint built in.
Here ${\cal N}_{i}(\gamma (t))$, $1\leq i\leq 2$, and ${\cal E}(\gamma (t))$
are density functionals determined by the solutions of the equations of
motion for the parameterizing functions ${\cal V}_{fib}(\gamma ,t)$ and $%
\omega (t)$. Appendix I gives their detailed structure.

The ${\cal N}_{i}(\gamma (t))$ in eq. (\ref{gammadot2}) are fiber quantities
that provide information on intra-fiber searches and their propagation in
time. Thus they act as time-dependent and density-dependent potentials. $%
{\cal E}(\gamma (t))$ also is an implicit function of fiber factors but an
explicit function of $\gamma $, hence may be decomposed in conventional KS
fashion into kinetic, coulombic, and exchange-correlation parts if that
proves useful.

\section{Summary and Conclusions}

We summarize what has been shown. First, to understand the details of
constrained search in the most general way, we established that the space of
trace class operators ${\cal B}_{1}({\cal H}^{N})$ can be sorted into
equivalence classes 
\begin{equation}
\left\{ \lbrack f]_{N};\quad f\in L_{1}({\tt Z})\right\}  \label{summequiv}
\end{equation}%
This classification endows ${\cal B}_{1}({\cal H}^{N})$ with two fiber
bundle structures, one affine ${\frak B}_{aff}$, the other vector ${\frak B}%
_{vect}$. In ${\frak B}_{aff}$, the fibers are the $[f]_{N}$ and the base of
the bundle is $L_{1}({\tt Z})$, while in ${\frak B}_{vect}$ the fibers are $%
\ker \Xi $ and the base is $\ker \Xi ^{c}$. The projection map from the
fiber to the base in the affine case is the contraction map $\Xi $ defined
in eq. (\ref{Xidefn}) 
\begin{equation}
\Xi (Y)=Tr\{\Phi ^{\dagger }\Phi Y\}=f;\quad Y\in {\cal B}_{1}({\cal H}%
^{N});\quad f\in L_{1}({\tt Z})  \label{summXidefn}
\end{equation}%
and the product of field operators is as in eq. (\ref{fieldopprod}). In the
vector bundle the corresponding projection map is the projection onto the
unique element $\Gamma \left( f\right) $ of $\ker \Xi ^{c}$ that corresponds
to $f$ \ under the mapping $\Xi $. The fibers of the affine bundle (the
equivalence classes $[f]_{N}$ ) correspond to $\left\{ \ker \Xi \right\}
\times \left\{ \Gamma \left( f\right) \right\} $ i.e. the set $\left\{
\left( k,\Gamma \left( f\right) \right) ;k\in \ker \Xi ,\text{ fixed }\Gamma
\left( f\right) \in \ker \Xi ^{c}\right\} $ which is a fiber and base point
of the vector bundle ${\frak B}_{vect}$.

Second, because $\Xi $ is a linear map, its kernel is a linear subspace of $%
{\cal B}_{1}({\cal H}^{N})$. Because of the fiber bundle structures, the
kernel plays a central role in executing the part of constrained search that
corresponds to searching at a specified density. Detailed knowledge of the
structure of the kernel, in particular its invariances, therefore provides
an entirely new route to construction of functionals.

Third, because physical systems are characterized by density operators, i.e.
positive trace class operators $Y\in {\cal B}_{1}({\cal H}^{N})_{+}\equiv
\left\{ Y\in {\cal B}_{1}({\cal H}^{N}),Y\geq 0\right\} $, we focused upon
fibers corresponding to $\left\{ [\gamma ]_{N},\;\gamma \geq 0\right\} $.
Then we demonstrated [Appendix E, eq. (\ref{decomp})] that, for pure states
(and all physical densities correspond to pure states, though not
necessarily ground states), any density operator can be decomposed as 
\begin{equation}
D^{N}=D_{K_{OD}}^{N}+D_{K_{D}}^{N}+\Gamma \left( \gamma \right)
\end{equation}

Fourth, each of these three contributions has specific invariance properties
related to whether the operators are diagonal or off-diagonal in the
coordinate eigenfunction representation and to whether they lie in $\ker \Xi 
$ or not; see Appendices E and F. For time-dependent DFT, these invariances
provide a Lie group structure that can be used to parameterize the dynamics
via the TDVP. For both time-dependent and time-independent DFT, the
invariances provide a new set of tools for constructing functionals. We
distinguish ``constructing'' from ``devising'' because Appendices G and I
comprise a {\em systematic} way to execute intra- and inter-fiber searches
via Lie groups of transformations. In both time-dependent and
time-independent forms of the theory, explicit knowledge of the optimal
result of the intra-fiber search is identical with determination of the
exact functional for the density which labels that particular fiber. When
restricted to a particular form of pure state, such knowledge is equivalent
to determination of the proper functional to associate with that type of
state.

The analytic form of those Lie group transformations is suggestive of
certain well-known cluster expansions.\cite{Blaizot86} This aspect is taken
up in Part II (following paper) as one of several means to obtain workable
systematic approximations. Also treated there are the related issues of the
various classical approximations for the nuclear motion (clamped nucleus,
adiabatic motion on electronic potential, narrow wave-packet, etc.). On that
topic we remark here only that an effort to extract one of those results
simply as the limiting case of the exact dynamical equations must be treated
with care because such a limit requires imposition of a constraint on the
Hamiltonian (or action), hence amounts to invoking a different Hamiltonian
to derive the dynamical equations.

Finally, the challenge of constructing physically realistic, yet tractable
time-dependent functionals is evident from the intricate structure of the $%
{\cal N}_{i} (\gamma (t))$ in eq. (\ref{gammadot2}). Further, because that
structure is related directly to the metric functions $\eta$, it is clear
that approximate functionals constructed on physical grounds or by
generalization from functionals used successfully in the time-independent
case actually must involve some implicit, unanalyzed assumptions about that
metric.

\section{Acknowledgements}

We thank Evaert Jan Baerends, Kieron Burke, Erik Deumens, Yngve \"{O}hrn,
A.K.~Rajagopal, and Robert van Leeuwen for informative and helpful
discussions. SBT was supported in part by the U.S. National Science
Foundation (DMR-KDI 9980015).

\appendix 

\section{Partitioned Equation of Motion in the ${\bf c}$ Parameterization}

Eq. (\ref{cgammaeom}), when rewritten in a form analogous with eq. (\ref%
{dEdlambda}), has the obvious matrix partitioning 
\begin{equation}
\left[ 
\begin{array}{cc}
\eta _{{\bf cc}}\left( t\right) & \eta _{{\bf c}\gamma }\left( t\right) \\ 
\eta _{\gamma {\bf c}}\left( t\right) & \eta _{\gamma \gamma }\left( t\right)%
\end{array}%
\right] \left[ 
\begin{array}{c}
{\bf \dot{c}} \\ 
\dot{\gamma}%
\end{array}%
\right] =\left[ 
\begin{array}{c}
\nabla _{{\bf c}}{\cal E}\left( t\right) \\ 
\nabla _{\gamma }{\cal E}\left( t\right)%
\end{array}%
\right]  \label{partmat1}
\end{equation}%
This is an operator notation and $\nabla _{\lambda }{\cal E}\left( t\right)
=\delta {\cal E}/\delta \lambda $. With $\lambda _{i}={\bf c}\,${\ or }$%
\gamma $, the components of the metric matrix $\eta $ are 
\begin{equation}
\eta _{\lambda _{1},\lambda _{2}}(t,{\bf y},{\bf y}^{\prime })=i\left\{ 
\frac{\delta }{\delta \tilde{\lambda}_{1}({\bf y})}\frac{\delta }{\delta
\lambda _{2}({\bf y}^{\prime })}-\frac{\delta }{\delta \lambda _{2}({\bf y})}%
\frac{\delta }{\delta \tilde{\lambda}_{1}({\bf y}^{\prime })}\right\} \ln S| 
\Sb \tilde{\lambda}_{1}=\lambda _{1}  \\ \tilde{\lambda}_{2}=\lambda _{2} 
\endSb   \label{etamat}
\end{equation}%
Here 
\begin{equation}
S\equiv \left\langle \Psi (\tilde{{\bf c}}(t),\tilde{\gamma}(t))|\Psi ({\bf c%
}(t),\gamma (t))\right\rangle
\end{equation}%
and the analogue to the classical Hamiltonian (recall eq. (\ref{Eoft})) is 
\begin{equation}
{\cal E}(t)=\frac{\left\langle \Psi ({\bf c}(t),\gamma (t))|H|\Psi ({\bf c}%
(t),\gamma (t))\right\rangle }{\left\langle \Psi ({\bf c}(t),\gamma
(t))|\Psi ({\bf c}(t),\gamma (t))\right\rangle }  \label{Eofcgammat}
\end{equation}%
Coupled equations of motion follow 
\begin{eqnarray}
\dot{{\bf c}}(t) &=&\eta _{{\bf c\,c}}^{-1}(t)\left[ \nabla _{{\bf c}}{\cal E%
}(t)-\eta _{{\bf c}\gamma }(t)\dot{\gamma}\right]  \nonumber \\
\dot{\gamma}(t) &=&\left[ \eta _{\gamma \gamma }(t)-\eta _{\gamma {\bf c}%
}(t)\eta _{{\bf cc}}^{-1}(t)\eta _{{\bf c}\,\gamma }(t)\right] ^{-1} 
\nonumber \\
&&\times \left[ \nabla _{\gamma }{\cal E}(t)-\eta _{\gamma {\bf c}}(t)\eta _{%
{\bf cc}}^{-1}(t)\nabla _{{\bf c}}{\cal E}(t)\right]  \label{cdotgammadot}
\end{eqnarray}%
Under the assumption that $\eta _{{\bf cc}}^{-1}(t)$ exists, denote by ${\bf %
c}_{\#}(t)$ the solution of the first of these for the correct density $%
\gamma _{\#}(t)$. Then the exact equation of motion for the temporal
development of the density is 
\begin{equation}
\dot{\gamma}(t)=\tilde{{\cal N}_{1}}(\gamma (t))\left[ \nabla _{\gamma }%
{\cal E}(\gamma (t))-\tilde{{\cal N}_{2}}(\gamma (t))\right]
\label{Appgammadot}
\end{equation}%
The functionals are defined as 
\begin{align}
\tilde{{\cal N}_{1}}(\gamma (t))& =\left[ \eta _{\gamma \gamma }({\bf c}%
_{\#}(t),\gamma (t))-\eta _{\gamma {\bf c}}({\bf c}_{\#}(t),\gamma (t))\eta
_{{\bf cc}}^{-1}({\bf c}_{\#}(t),\gamma (t))\right. \times  \nonumber \\
& \qquad \qquad \qquad \qquad \qquad \qquad \qquad \qquad \qquad \qquad
\left. \eta _{{\bf c}\gamma }({\bf c}_{\#}(t),\gamma (t))\right] ^{-1} 
\nonumber \\
\tilde{{\cal N}_{2}}(\gamma (t))& =\eta _{\gamma {\bf c}}({\bf c}%
_{\#}(t),\gamma (t))\eta _{{\bf cc}}^{-1}({\bf c}_{\#}(t),\gamma (t))\nabla
_{{\bf c}}{\cal E}({\bf c}_{\#}(t),\gamma (t)).\quad
\end{align}%
Any singularities are associated with bifurcation points of the time
evolution of the system.

\section{Unitary Groups and Invariances of the Kernel}

\noindent Here we give details as to why it is necessary to use the general
linear group rather than the unitary group to describe the invariances of
the kernel. The time-dependence is inconsequential to the discussion, so is
suppressed.

If consideration were restricted to unitary operators, the invariance group
of the kernel would be 
\begin{equation}
{\cal G}_{\ker }^{U}=\left\{ U\in {\cal U}({\cal H}^{N});{\rm Tr}\{\Phi
^{\dagger }\Phi UYU^{\dagger }\}={\rm Tr}\{\Phi ^{\dagger }\Phi
Y\}=0\,\forall Y\in \ker \Xi \right\}  \label{Gkerdefn}
\end{equation}%
where ${\cal U}({\cal H}^{N})$ is the group of unitary operators on $N$%
-particle space. The condition expressed in eq. $\left( \ref{Gkerdefn}%
\right) $ does not imply that transformations belonging to ${\cal G}_{\ker
}^{U}$ send elements of an equivalence class to itself 
\begin{equation}
{\cal G}_{\ker }^{U}:\left[ f\right] _{N}\rightarrow \left[ f^{\prime }%
\right] _{N};\quad \text{in general }f\neq f^{\prime }
\end{equation}%
On the other hand the condition that defines the unitary group 
\begin{equation}
{\cal G}_{\Xi }:\ker \Xi \rightarrow \ker \Xi
\end{equation}%
that does send equivalence classes to equivalence classes is given by 
\begin{equation}
{\cal G}_{\Xi }=\left\{ U\in {\cal U}({\cal H}^{N});{\rm Tr}\{\Phi ^{\dagger
}\Phi UYU^{\dagger }\}={\rm Tr}\{\Phi ^{\dagger }\Phi Y\}\quad \,\forall
Y\in {\cal B}_{1}\left( {\cal H}^{N}\right) \right\}  \label{GXi}
\end{equation}%
Evidently from eq. $\left( \ref{GXi}\right) $ 
\begin{equation}
\left[ U(\omega ),\Phi ^{\dagger }\Phi \right] =0
\end{equation}%
and ${\cal G}_{\Xi }$ is obviously a subgroup of ${\cal G}_{\ker }^{U}$ and
can be represented by transformations of the form 
\begin{eqnarray}
{\cal G}_{\Xi } &\subset &{\cal G}_{\ker }^{U}=\{U(\omega );\hat{U}(\omega
)(Y)\equiv U(\omega )YU^{\dagger }(\omega );  \nonumber \\
U(\omega ) &=&e^{i\,\int \omega (\kappa _{1},\ldots ,\kappa _{N})\Phi
^{\dagger }(\kappa _{1})\Phi (\kappa _{1})\ldots \Phi ^{\dagger }(\kappa
_{N})\Phi (\kappa _{N})d\kappa _{1}\ldots d\kappa _{N}};  \nonumber \\
&&\qquad \qquad \qquad \qquad \quad \omega (\kappa _{1},\ldots ,\kappa
_{N})\in {\Bbb R}\}  \label{GXidefn}
\end{eqnarray}%
One can see that ${\cal G}_{\Xi }$ is the largest commutative unitary
subgroup of ${\cal G}_{\ker }^{U}$ with these properties. Because of eq. (%
\ref{WifromYref}), it follows that fibers associated with states are mapped
into themselves by $\hat{U}(\omega )$ 
\begin{equation}
Y\in \lbrack \gamma ]_{N}\Leftrightarrow \hat{U}(\omega )Y\in \lbrack \gamma
]_{N}
\end{equation}%
However, given a reference element $Y_{ref}\in \lbrack \gamma ]_{N}$
corresponding to a pure state $|\Psi _{ref}\rangle $ in the fiber $[\gamma
]_{N}$, it is {\em not} the case that all the rest of the pure-state segment
of the fiber is generated by application of elements of ${\cal G}_{\Xi }$,
(see Appendix F), i.e. this group does {\em not }act homogeneously on the
set of pure states in $[\gamma ]_{N}$. This also shows that{\em \ there is
no intra fiber unitary subgroup of }${\cal G}_{\ker }^{U}$ {\em that acts
homogeneously on the set of pure states contained within a fiber}.{\em \ }%
The physical insight is that a transformation such as eq. (\ref{GXidefn})
never can generate the equivalent of a Jastrow or Hylleras factor in the
Hilbert space of pure states. If one drops the unitary requirement in eq. $%
\left( \ref{Gkerdefn}\right) $ {\it it is} possible find transformations,
that are diagonal in the coordinate representation (like those of ${\cal G}%
_{\Xi }$){\it , }which produce all pure states with a given density, $\gamma 
$, from a pure reference state $\left| \Psi \right\rangle $ of that density%
{\it . }These transformations are of the form%
\begin{equation}
T\left( {\cal V},\omega \right) =e^{\int \left\{ {\cal V}(\kappa _{1},\ldots
,\kappa _{N})+i\omega (\kappa _{1},\ldots ,\kappa _{N})\right\} \Phi
^{\dagger }(\kappa _{1})\Phi (\kappa _{1})\ldots \Phi ^{\dagger }(\kappa
_{N})\Phi (\kappa _{N})d\kappa _{1}\ldots d\kappa _{N}}
\end{equation}%
where there are certain restriction placed on ${\cal V}(\kappa _{1},\ldots
,\kappa _{N})$ (see Appendix F). We denote the group formed by these intra
fiber transformations ${\frak G}_{fib}$ in the main body of this article.
With a weaker set of restrictions, (see Appendix F), such maps also
transform $\ker \Xi \rightarrow \ker \Xi $, and form a subgroup ${\frak G}%
_{\ker }\subset {\frak G}_{\Xi }$, and with a third set of restrictions, one
forms a group of inter fiber transformations ${\frak G}_{den}$ (see Appendix
F).

\section{Positive Maps of Trace Class Operators}

\noindent Recall that a trace class operator $Y \; \in {\cal B}_{1} ( {\cal H%
})$, with ${\cal H}$ any Hilbert space, is positive $Y \ge 0$ if 
\begin{equation}
\langle v | Y v \rangle \ge 0 \qquad \forall \, v \in {\cal H}
\end{equation}
which is true if and only if all the eigenvalues of $Y \ge 0$. In turn this
requirement is true if and only if 
\begin{equation}
Y = Q Q^{\dagger} \, , Q\in {\cal B}_{2} ( {\cal H})
\end{equation}
and ${\cal B}_{2} ( {\cal H})$ is the space of Hilbert-Schmidt operators.

A linear map of trace class operators $\hat{T}:{\cal B}_{1}({\cal H})\mapsto 
{\cal B}_{1}({\cal H})$ is positive if $\hat{T}(Y)\geq 0$ for $Y\geq 0$. In
particular, $\hat{T}(Y)=TYT^{\dagger }$ is a positive map ${\cal B}_{1}(%
{\cal H})\mapsto {\cal B}_{1}({\cal H})$ for any bounded operator $T$, i.e.,
for such a $T$ 
\begin{equation}
\hat{T}(Y)\in {\cal B}_{1}({\cal H})_{+}\,\,\text{{\ if }}\,Y\in {\cal B}%
_{1}({\cal H})_{+}
\end{equation}%
where ${\cal B}_{1}({\cal H})_{+}$ is the cone of positive trace class
operators \cite{Kummer}. Cones facilitate analysis because they define order
relationships. That is, $A>B$ if $A-B\in {\cal B}_{1}({\cal H})_{+}$. The
set of positive operators is a cone, not just a convex set, as normalization
does not have to be constant.

If $T$ is an invertible bounded operator, i.e. $T\in GL({\cal H})$, then 
\begin{equation}
T=e^{A}\,e^{iB}
\end{equation}%
with $A^{\dagger }=A$ and $B^{\dagger }=B$. Clearly $e^{iB}$ is unitary
while $e^{A}$ is invertible and positive. The set $\{e^{iB}\}$ forms the
unitary subgroup ${\cal U}({\cal H})$ of $GL({\cal H})$, while the set $\{
e^{A}\}$ comprises the left coset $GL({\cal H})/{\cal U}({\cal H})$ of $%
{\cal U}({\cal H})$ in $GL({\cal H})$. Denote the maps $\hat{T}$ generated
by $T\in GL({\cal H})$ as $\widehat{GL({\cal H})}$ and observe that they map 
${\cal B}_{1}({\cal H})\mapsto {\cal B}_{1}({\cal H})$ and ${\cal B}_{2}(%
{\cal H})\mapsto {\cal B}_{2}({\cal H}).$ These are used in the next
Appendix.

\section{Group Orbits of Density Operators}

\noindent Once again the time-dependence appears only as a parameterization
that is inconsequential to the discussion, so is omitted. The transformation 
\begin{eqnarray}
\hat{T} &:&{\cal B}_{1}({\cal H})\mapsto {\cal B}_{1}\left( {\cal H}\right) 
\nonumber \\
&&\hat{T}(Y)=TYT^{\dagger }
\end{eqnarray}%
maintains the rank of the operator $Y$, that is $\hat{T}(Y)$ has the same
number of non-zero eigenvalues as $Y$ itself. Particularizing to density
operators $D$, their cone ${\cal B}_{1}({\cal H})_{+}$ can be classified by
the rank of $D\in {\cal B}_{1}({\cal H})_{+}$. Operators for pure states
have unit rank and correspond to vectors in ${\cal H}$ 
\begin{equation}
\mathop{\rm rank}\nolimits(D)=1\,\,\Leftrightarrow D=\frac{|\Psi \rangle
\langle \Psi |}{\langle \Psi |\Psi \rangle }\,;|\Psi \rangle \in {\cal H}
\end{equation}

All rank 1 density operators lie in the same orbit of $\widehat{GL({\cal H})}
$ (see Appendix C for definition) in ${\cal B}_{1}({\cal H})$. That is, if $%
D_{0}$ is an arbitrary, specified pure state density operator, then all
other pure state density operators $D$ are generated by 
\begin{equation}
D=\hat{T}(D_{0});\quad \hat{T}\in \widehat{GL({\cal H})}\,;\quad \mathop{\rm
rank}\nolimits (D_{0})=1
\end{equation}

Ensemble states correspond to density operators with rank greater than
unity. As in the pure state case, ensemble density operators of a given rank
all lie on one orbit of $\widehat{GL({\cal H})}$. Thus, if $D_{p,0}$ is an
arbitrary, specified ensemble state density operator of rank $p$, then all
other ensemble state density operators of the same rank $D_{p}$ are
generated by 
\begin{equation}
D_{p}=\hat{T}(D_{p,0});\quad \hat{T}\in \widehat{GL({\cal H})}\,;\quad %
\mathop{\rm rank}\nolimits (D_{p,0})=p
\end{equation}

\section{Analytical Structure of the Kernel}

\noindent This Appendix treats characteristics of $Y\in \ker \Xi $. As
before, the time-dependence appears only as a parameterization that is
inconsequential to the discussion, so is omitted. It is useful for present
purposes to begin by expressing all operators belonging to ${\cal B}_{1}(%
{\cal H}^{N})$ (recall $N=N_{e}+N_{n}$) in terms of the ordered generalized
operator basis \cite{BerezinShubin} 
\begin{equation}
\left\{ |{\bf x}_{1}\ldots {\bf x}_{N_{e}}\rangle _{ord}\langle {\bf x}%
_{1}^{\prime }\ldots {\bf x}_{N_{e}}^{\prime }|_{ord}\otimes |{\bf X}%
_{1}\ldots {\bf X}_{N_{n}}\rangle _{ord}\langle {\bf X}_{1}^{\prime }\ldots 
{\bf X}_{N_{n}}^{\prime }|_{ord}\right\}
\end{equation}%
with 
\begin{equation}
{\bf x},{\bf x}^{\prime },{\bf X},{\bf X}^{\prime }\in {\Bbb F}
\end{equation}%
where ${\Bbb F}={\Bbb R}^{3}\otimes {\Bbb C}^{2}$ if the space-spin analysis
is pursued and ${\Bbb F}={\Bbb R}^{3}$ if spin is suppressed (whence the
states are coordinate eigenfunctions). The subscript ``$ord$'' denotes the
ordered basis 
\begin{equation}
{\bf x}_{1}\leq {\bf x}_{2}\leq {\bf x}_{3}\leq \ldots
\end{equation}%
with the understanding that strict inequality holds only for each species of
fermions. Shortly a shift to the unordered basis will be advantageous. For
either basis, a compressed notation, analogous with eq.~(\ref{compressed}),
is convenient 
\begin{eqnarray}
&&|{\bf x}_{1}\ldots {\bf x}_{N_{e}}\rangle \langle {\bf x}_{1}^{\prime
}\ldots {\bf x}_{N_{e}}^{\prime }|\otimes |{\bf X}_{1}\ldots {\bf X}%
_{N_{n}}\rangle \langle {\bf X}_{1}^{\prime }\ldots {\bf X}_{N_{n}}^{\prime
}|  \nonumber \\
&\equiv &|\kappa _{1}\ldots \kappa _{N_{e}},\kappa _{N_{e}+1}\ldots \kappa
_{N}\rangle \langle \kappa _{1}^{\prime }\ldots \kappa _{N_{e}}^{\prime
},\kappa _{N_{e}+1}^{\prime }\ldots \kappa _{N}^{\prime }|  \nonumber \\
&\equiv &|{\sf x},{\sf X}\rangle \langle {\sf x}^{\prime },{\sf X}^{\prime
}|,
\end{eqnarray}%
The last form (recall eq. (\ref{manifdef})) will be used only when there is
no ambiguity.

Any operator $Y\,\in {\cal B}_{1}({\cal H}^{N})$ can be expanded in terms of
the ordered basis as 
\begin{eqnarray}
Y &=&\int \,d\kappa _{1},\ldots ,d\kappa _{N},d\kappa _{1}^{\prime },\ldots
,d\kappa _{N}^{\prime }\tilde{Y}(\kappa _{1},\ldots ,\kappa _{N};\kappa
_{1}^{\prime },\ldots ,\kappa _{N}^{\prime })  \nonumber \\
&&\times |\kappa _{1}\ldots \kappa _{N}\rangle _{ord}\langle \kappa
_{1}^{\prime }\ldots \kappa _{N}^{\prime }|_{ord}
\end{eqnarray}%
The operator integral kernel with respect to the ordered basis can be
expanded \cite{Richtmyer} as 
\begin{eqnarray}
\tilde{Y}(\kappa _{1},\ldots ,\kappa _{N},\kappa _{1}^{\prime },\ldots
,\kappa _{N}^{\prime }) &=&\sum_{i=1}^{\infty }\alpha _{i}\bar{\varphi}%
_{i}(\kappa _{1},\ldots ,\kappa _{N})\psi _{i}(\kappa _{1}^{\prime },\ldots
,\kappa _{N}^{\prime });  \nonumber \\
\alpha _{i} &\geq &0; \qquad \sum_{i=1}^{\infty }\alpha _{i}<\infty
\end{eqnarray}%
where the basis functions $\left\{ \varphi _{i},\psi _{i}\right\} $ are
specific to the operator $Y$ in the sense of spanning its range and domain
and satisfying 
\begin{eqnarray}
\int \,d\kappa _{1},\ldots ,d\kappa _{N},\bar{\varphi}_{i}(\kappa
_{1},\ldots ,\kappa _{N})\varphi _{j}(\kappa _{1},\ldots ,\kappa _{N})
&=&\delta _{ij}  \nonumber \\
\int \,d\kappa _{1},\ldots ,d\kappa _{N},\bar{\psi}_{i}(\kappa _{1},\ldots
,\kappa _{N})\psi _{j}(\kappa _{1},\ldots ,\kappa _{N}) &=&\delta _{ij}
\end{eqnarray}

It is easier to consider the contraction map, the central quantity, in the 
{\em unordered} basis. In it, the operator integral kernel is 
\begin{equation}
Y(\kappa _{1},\ldots ,\kappa _{N},\kappa _{1}^{\prime },\ldots ,\kappa
_{N}^{\prime })=\frac{(-1)^{(P_{e}+P_{n})}}{(N_{e}!N_{n}!)^{2}}{\cal P}%
\tilde{Y}(\kappa _{1},\ldots ,\kappa _{N},\kappa _{1}^{\prime },\ldots
,\kappa _{N}^{\prime })
\end{equation}%
where ${\cal P}^{-1}$ is the permutation operator that generates the ordered
strings from the unordered ones 
\begin{eqnarray}
{\cal P}^{-1}\{\kappa _{1},\ldots ,\kappa _{N}\} &=&\{\kappa _{1},\ldots
,\kappa _{N}\}_{ord}  \nonumber \\
{\cal P}^{-1}\{\kappa _{1}^{\prime },\ldots ,\kappa _{N}^{\prime }\}
&=&\{\kappa _{1}^{\prime },\ldots ,\kappa _{N}^{\prime }\}_{ord}
\end{eqnarray}%
$P_{e}$ is the number of permutations of the electron coordinates, and $%
P_{n} $ is either the number of permutations of the nuclear coordinates if
the nuclei are fermions or $0$ if they are bosons.

In the unordered basis the contraction map $\Xi $ eq. (\ref{Xidefn}) on any
operator $Y$ is 
\begin{align}
& \!\!\!\!\!Tr\{\Phi ^{\dagger }({\bf x})\Phi ({\bf x})\Phi ^{\dagger }({\bf %
X})\Phi ({\bf X})Y\}  \nonumber \\
& =\int \,d\kappa _{1},\ldots ,d\kappa _{N}d\kappa _{1}^{\prime },\ldots
,d\kappa _{N}^{\prime }d\kappa _{1}^{\prime \prime },\ldots ,d\kappa
_{N}^{\prime \prime }  \nonumber \\
& \qquad \times \langle \kappa _{1}^{\prime \prime },\ldots ,\kappa
_{N}^{\prime \prime }|\Phi ^{\dagger }({\bf x})\Phi ({\bf x})\Phi ^{\dagger
}({\bf X})\Phi ({\bf X})|\kappa _{1},\ldots ,\kappa _{N}\rangle  \nonumber \\
& \qquad \times \langle \kappa _{1}^{\prime },\ldots ,\kappa _{N}^{\prime
}|\kappa _{1}^{\prime \prime },\ldots ,\kappa _{N}^{\prime \prime }\rangle
\;Y(\kappa _{1},\ldots ,\kappa _{N};\kappa _{1}^{\prime },\ldots ,\kappa
_{N}^{\prime })  \nonumber \\
& =N_{e}N_{n}\,\int d{\bf x}_{2}^{\prime }\ldots ,d{\bf x}_{N_{e}}^{\prime }d%
{\bf X}_{2}^{\prime },\ldots ,d{\bf X}_{N_{n}}^{\prime }  \nonumber \\
& \times Y({\bf x},{\bf x}_{2}\ldots ,{\bf x}_{N_{e}}{\bf X},{\bf X}%
_{2}\ldots ,{\bf X}_{N_{n}};{\bf x},{\bf x}_{2}^{\prime }\ldots ,{\bf x}%
_{N_{e}}^{\prime }{\bf X},{\bf X}_{2}^{\prime }\ldots ,{\bf X}%
_{N_{n}}^{\prime })  \label{TrphiphiphiphiY}
\end{align}

The elements of $\ker \Xi $ become evident by writing 
\begin{align}
Y(\kappa _{1},\ldots ,\kappa _{N};\kappa _{1}^{\prime },\ldots ,\kappa
_{N}^{\prime })& =Y_{d}(\kappa _{1},\ldots ,\kappa
_{N})\prod_{j=1}^{N}\delta (\kappa _{j}-\kappa _{j}^{\prime })  \nonumber \\
& +Y_{OD}(\kappa _{1},\ldots ,\kappa _{N};\kappa _{1}^{\prime },\ldots
,\kappa _{N}^{\prime })  \label{YYdYod}
\end{align}%
\begin{eqnarray}
Y_{OD}(\kappa _{1},\ldots ,\kappa _{N};\kappa _{1},\ldots ,\kappa _{N}) &=&0
\nonumber \\
Y_{d}(\kappa _{1},\ldots ,\kappa _{N}) &\equiv &Y(\kappa _{1},\ldots ,\kappa
_{N};\kappa _{1},\ldots ,\kappa _{N})
\end{eqnarray}%
and the caveat that the product of delta functions in eq. (\ref{YYdYod})
implicitly includes all the permutations of coordinates (because of working
in the unordered basis). Note that $Y_{d}$ is symmetric with respect to
interchange among the electron coordinates themselves and similarly for the
nuclear coordinates. Substitution of (\ref{YYdYod}) in (\ref{TrphiphiphiphiY}%
) gives 
\begin{gather}
Tr\!\{\Phi ^{\dagger }({\bf x})\Phi ({\bf x})\Phi ^{\dagger }({\bf X})\Phi (%
{\bf X})Y\}=  \nonumber \\
N_{e}N_{n}\,\int d{\bf x}_{2}^{\prime }\ldots ,d{\bf x}_{N_{e}}^{\prime }d%
{\bf X}_{2}^{\prime },\ldots ,d{\bf X}_{N_{n}}^{\prime }Y_{d}({\bf x},{\bf x}%
_{2}\ldots ,{\bf x}_{N_{e}}{\bf X},{\bf X}_{2}\ldots ,{\bf X}_{N_{n}})
\label{ContrYtoYd}
\end{gather}%
An alternative way to state this is that the contraction map sends all
off-diagonal basis operators to zero 
\begin{eqnarray}
&&\Xi \left\{ |\kappa _{1}\ldots \kappa _{N}\rangle \langle \kappa
_{1}^{\prime }\ldots \kappa _{N}^{\prime }|\right\} =0  \nonumber \\
&&\;\text{{for }}\left\{ \kappa _{1}\ldots \kappa _{N_{e}}\right\} \neq
\left\{ \kappa _{1}^{\prime }\ldots \kappa _{N_{e}}^{\prime }\right\} \text{
or }\left\{ \kappa _{1}\ldots \kappa _{N_{n}}\right\} \neq \left\{ \kappa
_{1}^{\prime }\ldots \kappa _{N_{n}}^{\prime }\right\}  \label{zeroeta}
\end{eqnarray}

Three cases are obvious immediately from eq. (\ref{ContrYtoYd}); two belong
to $\ker \Xi $, one does not. First, there are pure off-diagonal operators 
\begin{equation}
Y=Y_{OD}\,\,;\,\,(Y_{d}=0)\,\Rightarrow Y\in \ker \Xi  \label{Ycase1}
\end{equation}%
Then there are diagonal operators with a special property 
\begin{gather}
Y=Y_{d};\,\,(Y_{OD}=0)  \nonumber \\
\int d{\bf x}_{2}\ldots ,d{\bf x}_{N_{e}}d{\bf X}_{2},\ldots ,d{\bf X}%
_{N_{n}}Y_{d}({\bf x},{\bf x}_{2}\ldots ,{\bf x}_{N_{e}}{\bf X},{\bf X}%
_{2}\ldots ,{\bf X}_{N_{n}})=0  \nonumber \\
\Rightarrow Y\in \ker \Xi  \label{Ycase2}
\end{gather}%
Finally, there are diagonal operators that do not belong to $\ker \Xi $ 
\begin{gather}
Y=Y_{d};\,\,(Y_{OD}=0)  \nonumber \\
0\neq \int d{\bf x}_{2}\ldots ,d{\bf x}_{N_{e}}d{\bf X}_{2},\ldots ,d{\bf X}%
_{N_{n}}Y_{d}({\bf x},{\bf x}_{2}\ldots ,{\bf x}_{N_{e}}{\bf X},{\bf X}%
_{2}\ldots ,{\bf X}_{N_{n}})  \nonumber \\
\Rightarrow Y\notin \ker \Xi  \label{Ycase3}
\end{gather}

The space $K_{OD}$ generated by operators satisfying (\ref{Ycase1}) is a
subspace of $\ker \{\Xi \}$ 
\begin{eqnarray}
K_{OD} &\subset &\ker \{\Xi \}  \nonumber \\
K_{OD} &=&{\rm linear\,span}\left\{ Y=Y_{OD},Y_{d}=0\right\}  \label{koddefn}
\end{eqnarray}%
The remaining part $K_{D}$ of $\ker \{\Xi \}$ is comprised of diagonal
operators that satisfy eq. (\ref{Ycase2}) 
\begin{eqnarray}
K_{D} &\subset &\ker \{\Xi \};\text{ where }K_{D}={\rm linear\,span}\left\{
Y=Y_{d},Y_{OD}=0;\right.  \nonumber \\
&&\left. 0=\int d{\bf x}_{2},\ldots d{\bf x}_{N_{e}}d{\bf X}_{2},\ldots d%
{\bf X}_{N_{n}}Y_{d}({\bf x},{\bf x}_{2},\ldots {\bf x}_{N_{e}}{\bf X},{\bf X%
}_{2},\ldots {\bf X}_{N_{n}})\right\}  \nonumber \\
&&  \label{kddefn}
\end{eqnarray}%
In detail, for $K_{D}$, $Y_{d}$ obeys at least one of the two conditions 
\begin{eqnarray}
0 &=&\int \,d{\bf x}_{i_{1}},\ldots ,d{\bf x}_{i_{N_{e}-1}},Y_{d}({\bf x}%
_{1},\ldots ,{\bf x}_{N_{e}}{\bf X}_{1},\ldots ,{\bf X}_{N_{n}})  \nonumber
\\
0 &=&\int \,d{\bf X}_{j_{1}},\ldots ,d{\bf X}_{j_{N_{n}-1}},Y_{d}({\bf x}%
_{1},\ldots ,{\bf x}_{N_{e}}{\bf X}_{1},\ldots ,{\bf X}_{N_{n}})
\label{kddefn1}
\end{eqnarray}%
with 
\begin{eqnarray}
\left\{ i_{1}\ldots i_{N_{e}-1}\right\} &\subset &\left\{ 1\ldots
N_{e}\right\}  \nonumber \\
\left\{ j_{1}\ldots j_{N_{n}-1}\right\} &\subset &\left\{ 1\ldots
N_{n}\right\}  \label{indices}
\end{eqnarray}

Thus the kernel of $\Xi $ can be expressed as 
\begin{equation}
\ker \{\Xi \}=K_{OD}\oplus K_{D}
\end{equation}%
with the two components orthogonal with respect to the Hilbert-Schmidt inner
product 
\begin{equation}
Tr\{k_{OD}^{\dagger }k_{D}\}=0\;{\rm or}\,\,k_{OD}\in K_{OD},\,\,k_{d}\in
K_{D}
\end{equation}%
This result follows from 
\begin{eqnarray}
Tr\{|\kappa _{1}\ldots \kappa _{N}\rangle \langle \kappa _{1}\ldots \kappa
_{N}||\kappa _{1}^{\prime }\ldots \kappa _{N}^{\prime }\rangle \langle
\kappa _{1}^{\prime \prime }\ldots \kappa _{N}^{\prime \prime }|\} &=& 
\nonumber \\
\langle \kappa _{1}\ldots \kappa _{N}|\kappa _{1}^{\prime }\ldots \kappa
_{N}^{\prime }\rangle \langle \kappa _{1}^{\prime \prime }\ldots \kappa
_{N}^{\prime \prime }|\kappa _{1}\ldots \kappa _{N}\rangle &=&0
\end{eqnarray}%
which holds when both 
\begin{equation}
\{\kappa _{1}^{\prime }\ldots \kappa _{N_{e}}^{\prime }\}\neq \{\kappa
_{1}^{\prime \prime }\ldots \kappa _{N_{e}}^{\prime \prime }\}
\end{equation}%
{\em and} 
\begin{equation}
\{\kappa _{N_{e}+1}^{\prime }\ldots \kappa _{N}^{\prime }\}\neq \{\kappa
_{N_{e}+1}^{\prime \prime }\ldots \kappa _{N}^{\prime \prime }\}
\end{equation}%
In either case the primed and double primed sets of coordinates have at
least one difference, hence both cannot be equal to $\{\kappa _{1}\ldots
\kappa _{N}\}$.

Eq. (\ref{Ycase3}) makes clear that the only operators $Y\in {\cal B}_{1}(%
{\cal H}^{N})$ that do not lie in $\ker \Xi $ are diagonal operators that do
not belong to $K_{D}$. These form the subspace (of ${\cal B}_{1}({\cal H}%
^{N})$) 
\begin{eqnarray}
S_{D} &=&{\rm linear\,span}\big\{Y=Y_{d},Y_{OD}=0;\big.\qquad  \nonumber \\
&&\big.0\neq \!\!\!\int d{\bf x}_{2}\ldots ,d{\bf x}_{N_{e}}d{\bf X}%
_{2},\ldots ,d{\bf X}_{N_{n}}Y_{d}({\bf x},{\bf x}_{2}\ldots ,{\bf x}_{N_{e}}%
{\bf X},{\bf X}_{2}\ldots ,{\bf X}_{N_{n}})\!\big\}  \nonumber \\
&&  \label{sddefn}
\end{eqnarray}%
Therefore the space of trace class operators has the direct sum
decomposition 
\begin{equation}
{\cal B}_{1}({\cal H}^{N})=K_{OD}\oplus K_{D}\oplus S_{D}  \label{decomp}
\end{equation}

\section{Vector Bundle Transformations}

\noindent This Appendix develops pertinent properties of transformations $%
TYT^{\dagger }$ of trace class operators $Y$ in $K_{OD}\oplus K_{D}$ $\oplus 
$ $S_{D}$ by elements $T$ of the General Linear Group $GL\left( {\cal H}%
^{N}\right) $ of ${\cal H}^{N}$. In the coordinate eigenfunction basis, the
transformation $TYT^{\dagger }$ reads 
\begin{equation}
TYT^{\dagger }\equiv \hat{T}Y=\int d{\sf z}d{\sf z}^{\prime }\left[
\,TYT^{\dagger }({\sf z},{\sf z}^{\prime })\,\right] |{\sf z}\rangle \langle 
{\sf z}^{\prime }|  \label{TYT}
\end{equation}%
where the transformed integral kernel is 
\begin{equation}
TYT^{\dagger }({\sf z},{\sf z}^{\prime })=\int d{\sf z}^{\prime \prime }d%
{\sf z}^{\prime \prime \prime }T({\sf z},{\sf z}^{\prime \prime })Y({\sf z}%
^{\prime \prime },{\sf z}^{\prime \prime \prime })T^{\ast }({\sf z}^{\prime
\prime \prime },{\sf z}^{\prime })  \label{TYTkernel}
\end{equation}%
Clearly the only transforms that can preserve the structure described by
eqs. (\ref{koddefn}), (\ref{kddefn}), (\ref{sddefn}) have diagonal
representations 
\begin{eqnarray}
T({\sf z},{\sf z}^{\prime }) &=&\delta ({\sf z}-{\sf z}^{\prime })\,T({\sf z}%
)  \nonumber \\
T({\sf z}) &=&T({\sf x},{\sf X})  \label{Tddefn}
\end{eqnarray}%
with $T({\sf z})$ a species symmetric function with respect to interchange
among electronic coordinates and similarly for nuclear coordinates (within
each species).

The explicit form of the diagonal operator is 
\begin{eqnarray}
T &=&e^{\int {\cal V}({\bf x}_{1},\ldots ,{\bf x}_{N_{e}}{\bf X}_{1},\ldots ,%
{\bf X}_{N_{n}})\prod\limits_{j=1}^{N_{e}}\Phi _{e}^{\dagger }(x_{j})\Phi
_{e}(x_{j})\prod\limits_{k=1}^{N_{n}}\Phi _{n}^{\dagger }(X_{k})\Phi
_{n}(X_{k})\prod\limits_{l=1}^{N_{e}}dx_{l}\prod_{m=1}^{N_{n}}dX_{m}} 
\nonumber \\
&&\;\times e^{i\int \omega ({\bf x}_{1},\ldots ,{\bf x}_{N_{e}}{\bf X}%
_{1},\ldots ,{\bf X}_{N_{n}})\prod\limits_{q=1}^{N_{e}}\Phi _{e}^{\dagger
}(x_{q})\Phi _{e}(x_{q})\prod\limits_{r=1}^{N_{n}}\Phi _{n}^{\dagger
}(X_{r})\Phi
_{n}(X_{r})\prod\limits_{s=1}^{N_{e}}dx_{s}\prod_{t=1}^{N_{n}}dX_{t}}
\label{Tdexplicit}
\end{eqnarray}%
When $T$ operates upon an element of the basis of many-body coordinate
eigenfunctions, the result is 
\begin{eqnarray}
&&T\left| y_{1}\ldots y_{N_{e}}Y_{1}\ldots Y_{N_{n}}\right\rangle  \nonumber
\\
&&\quad =e^{{\cal V}(y_{1}\ldots y_{N_{e}}Y_{1}\ldots Y_{N_{n}})}e^{i\omega
(y_{1}\ldots y_{N_{e}}Y_{1}\ldots Y_{N_{n}})}\left| y_{1}\ldots
y_{N_{e}}Y_{1}\ldots Y_{N_{n}}\right\rangle  \label{Tdcoord}
\end{eqnarray}%
because 
\begin{eqnarray}
\Phi _{e}^{\dagger }(x_{j})\Phi _{e}(x_{j})\left| y_{1}\ldots
y_{N_{e}}Y_{1}\ldots Y_{N_{n}}\right\rangle &=&\left| y_{1}\ldots
y_{N_{e}}Y_{1}\ldots Y_{N_{n}}\right\rangle  \nonumber \\
\qquad \qquad \qquad \qquad \qquad \qquad \qquad \qquad \text{{for}}\;x_{j}
&\in &\{y_{1}\ldots y_{N_{e}}\}  \nonumber \\
&=&0\,\,\text{o{therwise}}
\end{eqnarray}%
and the corresponding relationship for $\Phi _{n}^{\dagger }(X_{k})\Phi
_{n}(X_{k})$.

For $Y=Y_{K_{OD}}\,\in K_{OD}$, diagonal transformations $\hat{T}$ map $%
K_{OD}\mapsto K_{OD}$ with the explicit form of the transformed integral
kernel 
\begin{equation}
TY_{K_{OD}}T^{\dagger }({\sf z},{\sf z}^{\prime })=T({\sf z})Y_{K_{OD}}({\sf %
z},{\sf z}^{\prime })T^{\ast }({\sf z}^{\prime })\in K_{OD}
\label{TYodTkernel}
\end{equation}%
Diagonal transformations that map $K_{D}\mapsto K_{D}$ provided that the
transformed integral kernel satisfies eq. (\ref{Ycase2}), which is in fact
the defining property for maps that maintain the vector bundle structure of $%
{\frak B}_{vec}$. This can be written as 
\begin{eqnarray}
\hat{T} &:&Y_{K_{D}}\mapsto TY_{K_{D}}T^{\dagger };\quad Y_{K_{D}}\in K_{D} 
\nonumber \\
&&TY_{K_{D}}T^{\dagger }({\sf z})=|T({\sf z})|^{2}Y_{K_{D}}({\sf z})
\end{eqnarray}%
{\em provided} that 
\begin{equation}
|T({\sf z})|^{2}=1+\phi ({\sf z})  \label{phidef}
\end{equation}%
and the species symmetric functions $\phi $ have the property 
\begin{align}
0& =\int \,d{\bf x}_{i_{1}},\ldots ,d{\bf x}_{i_{N_{e}-1}},d{\bf X}%
_{j_{1}},\ldots ,d{\bf X}_{j_{N_{n}-1}}  \nonumber \\
& \times \phi ({\bf x}_{1},\ldots ,{\bf x}_{N_{e}}{\bf X}_{1},\ldots ,{\bf X}%
_{N_{n}})Y_{d}({\bf x}_{1},\ldots ,{\bf x}_{N_{e}}{\bf X}_{1},\ldots ,{\bf X}%
_{N_{n}})  \nonumber \\
& \qquad \qquad \qquad \qquad \qquad \qquad \qquad \qquad \qquad \quad
\forall Y_{K_{D}}\in K_{D}  \label{strongorthog}
\end{align}%
and the indices on the variables of integration obey eq. (\ref{indices}).
These conditions follow directly from eqs. (\ref{Ycase2}). and (\ref%
{strongorthog}) and expresses the property of strong orthogonality, denoted
by $\breve{\phi}\perp _{S}K_{D}$, or equivalently, 
\begin{equation}
\breve{\phi}\perp _{S}\forall Y_{K_{D}}\in K_{D}
\end{equation}%
where $\breve{\phi}$ is the operator defined by the integral kernel $\phi $.
In particular, eq. (\ref{strongorthog}) implies that the pointwise product
function $\phi Y_{d}\in L_{1}({\tt Z})$, which in turn implies that $\phi
\in L_{\infty }({\tt Z})$ since $Y({\sf z})\in L_{1}({\tt Z})$. Here 
\begin{equation}
{\tt Z}={\Bbb R}^{3N_{e}}\times {\Bbb R}^{3N_{n}}\times {\Bbb C}%
^{2N_{e}}\times {\Bbb C}^{2N_{n}}
\end{equation}

Vector bundle transformations built from $\phi $ satisfy 
\begin{gather}
\phi \in L_{\infty }({\tt Z})  \nonumber \\
\phi \perp _{S}K_{D}  \label{phiprop}
\end{gather}%
and have the property of mapping $K_{D}\mapsto K_{D}$, $S_{D}\mapsto
K_{D}\oplus S_{D}$ and form the group ${\cal G}_{\ker }$. This group
together with the unitary group ${\cal G}_{\Xi }$ forms the group ${\frak G}%
_{\ker }={\cal G}_{\ker }\times {\cal G}_{\Xi }$, which is a subgroup of the
commutative subgroup ${\cal C}({\cal H}^{N})\subset GL({\cal H}^{N})$. The
elements of ${\cal C}({\cal H}^{N})$ are given by 
\begin{equation}
T=e^{\left\{ {\cal \breve{V}}+i\breve{\omega}\right\} }=T_{d}e^{i\breve{%
\omega}}
\end{equation}%
where the operator $T_{d}$ is defined by the integral kernel 
\begin{equation}
T_{d}(\kappa _{1}\ldots \kappa _{N})=\exp [{\cal V}(\kappa _{1}\ldots \kappa
_{N})]=(1+\phi (\kappa _{1}\ldots \kappa _{N}))^{\frac{1}{2}}
\end{equation}%
with no restrictions on $\phi $. The integral kernel of ${\cal \breve{V}}$
is given by 
\begin{equation}
{\cal V}(\kappa _{1},\ldots ,\kappa _{N})\equiv \frac{1}{2}\ln (1+\phi
(\kappa _{1},\ldots ,\kappa _{N}))  \label{Vhalfln}
\end{equation}%
The transformations $\hat{T}:K_{OD}\mapsto K_{OD}$ are determined by
elements of ${\frak G}_{\ker }$ and act according to eq. $\left( \ref%
{TYodTkernel}\right) $.

We now develop the conditions $\phi $ must satisfy to determine the
subgroups ${\frak G}_{den}$ and ${\frak G}_{fib}$ of inter-fiber and
intra-fiber maps respectively contained in ${\frak G}_{\ker }$. Once this is
accomplished we can construct a non redundant parameterizations of pure
states in the equivalence classes $\left[ \gamma \right] _{N}$ which can be
utilized in the constrained search procedure.

First consider transformations, $\hat{T}_{den}\left( \phi \right) $, between
fibers that produce the subgroup ${\frak G}_{den}$, eqs. (\ref{phidef}), (%
\ref{phiprop}) give the conditions for $\phi (\kappa _{1},\ldots ,\kappa
_{N})$ to map $S_{D}\mapsto S_{D}$. In that case 
\begin{eqnarray}
&&\int \phi ({\bf x}_{1},\ldots ,{\bf x}_{N_{e}}{\bf X}_{1},\ldots ,{\bf X}%
_{N_{n}})Y_{d}({\bf x}_{1},\ldots ,{\bf x}_{N_{e}}{\bf X}_{1},\ldots ,{\bf X}%
_{N_{n}})  \nonumber \\
&&\qquad \qquad \qquad \qquad \times d{\bf x}_{2},\ldots ,d{\bf x}_{N_{e}},d%
{\bf X}_{2},\ldots ,d{\bf X}_{N_{n}}  \nonumber \\
&=&\int \phi ({\bf x}_{1},\ldots ,{\bf x}_{N_{e}}{\bf X}_{1},\ldots ,{\bf X}%
_{N_{n}})Y_{S_{D}}({\bf x}_{1},\ldots ,{\bf x}_{N_{e}}{\bf X}_{1},\ldots ,%
{\bf X}_{N_{n}})  \nonumber \\
&&\qquad \qquad \qquad \qquad \times d{\bf x}_{2},\ldots ,d{\bf x}_{N_{e}},d%
{\bf X}_{2},\ldots ,d{\bf X}_{N_{n}}\neq 0
\end{eqnarray}%
with%
\begin{equation}
Y_{d}=Y_{K_{D}}+Y_{S_{D}}  \label{YKdYSd}
\end{equation}%
This corresponds to 
\begin{equation}
\Phi _{S_{D}}\left( {\sf z}_{1},{\sf z}_{2}\right) =\int P_{S_{D}}\left( 
{\sf z}_{1},{\sf z}_{1}^{\prime }\right) \phi \left( {\sf z}_{1}^{\prime
}\right) P_{S_{D}}\left( {\sf z}_{1}^{\prime },{\sf z}_{2}\right) d{\sf z}%
_{1}^{\prime }\neq 0\quad \forall {\sf z}_{1},{\sf z}_{2}  \label{projS}
\end{equation}%
and hence $\phi $ is not strongly orthogonal to $Y_{d}$, where $P_{S_{D}}$
is the integral kernel of the projector onto the subspace $S_{D}$. In order
to obtain the defining representation of ${\frak G}_{den}$ on $K_{D}\oplus
S_{D}$, the functions $\phi $ should have zero effect on $K_{D}$ thus need
to satisfy 
\begin{equation}
\Phi _{K_{D}}\left( {\sf z}_{1},{\sf z}_{2}\right) =\int P_{K_{D}}\left( 
{\sf z}_{1},{\sf z}_{1}^{\prime }\right) \phi \left( {\sf z}_{1}^{\prime
}\right) P_{K_{D}}\left( {\sf z}_{1}^{\prime },{\sf z}_{2}\right) d{\sf z}%
_{1}^{\prime }=0\quad \forall {\sf z}_{1},{\sf z}_{2}  \label{projK}
\end{equation}%
and 
\begin{equation}
\Phi _{K_{D}S_{D}}\left( {\sf z}_{1},{\sf z}_{2}\right) =\int
P_{K_{D}}\left( {\sf z}_{1},{\sf z}_{1}^{\prime }\right) \phi \left( {\sf z}%
_{1}^{\prime }\right) P_{S_{D}}\left( {\sf z}_{1}^{\prime },{\sf z}%
_{2}\right) d{\sf z}_{1}^{\prime }=0\quad \forall {\sf z}_{1},{\sf z}_{2}
\label{projKS}
\end{equation}%
The effect of these transformations on the off diagonal part of ${\cal B}%
_{1}\left( {\cal H}^{N}\right) $, are given by elements of ${\cal G}_{\Xi }$
that act according to eq. $\left( \ref{TYodTkernel}\right) $ and is given by 
\begin{equation}
\hat{T}_{den}\left( \phi \right) Y_{K_{OD}}=e^{{\cal \breve{V}}_{den}\left(
\phi \right) }Y_{K_{OD}}e^{{\cal \breve{V}}_{den}\left( \phi \right) }
\end{equation}

In order to produce intra-fiber transformations, $\hat{T}_{fib}\left( \phi
\right) :K_{D}\mapsto K_{D}$, that form the subgroup ${\cal G}_{fib}$ the
functions $\phi $ must satisfy the condition 
\begin{equation}
\Phi _{S_{D}}\left( {\sf z}_{1},{\sf z}_{2}\right) =\int P_{S_{D}}\left( 
{\sf z}_{1},{\sf z}_{1}^{\prime }\right) \phi \left( {\sf z}_{1}^{\prime
}\right) P_{S_{D}}\left( {\sf z}_{1}^{\prime },{\sf z}_{2}\right) d{\sf z}%
_{1}^{\prime }=0\quad \forall {\sf z}_{1},{\sf z}_{2}
\end{equation}%
The defining representation of ${\cal G}_{fib}$ on $K_{D}\oplus S_{D}$, is
obtained when the functions further satisfy 
\begin{equation}
\Phi _{K_{D}S_{D}}\left( {\sf z}_{1},{\sf z}_{2}\right) =\int
P_{K_{D}}\left( {\sf z}_{1},{\sf z}_{1}^{\prime }\right) \phi \left( {\sf z}%
_{1}^{\prime }\right) P_{S_{D}}\left( {\sf z}_{1}^{\prime },{\sf z}%
_{2}\right) d{\sf z}_{1}^{\prime }=0\quad \forall {\sf z}_{1},{\sf z}_{2}
\end{equation}%
The effect of these transformations on the off diagonal subspace $K_{OD}$
and is given by 
\begin{equation}
\hat{T}_{fib}\left( \phi \right) Y_{K_{OD}}=e^{{\cal \breve{V}}_{fib}\left(
\phi \right) }Y_{K_{OD}}e^{{\cal \breve{V}}_{fib}\left( \phi \right) }
\end{equation}%
The group of intra fiber transformations that acts homogeneously on the set
of pure states within a fiber is ${\frak G}_{fib}={\cal G}_{fib}\times {\cal %
G}_{\Xi }$ i.e is constructed from transformations of the form $e^{{\cal 
\breve{V}}_{fib}\left( \phi \right) }e^{i\breve{\omega}}$

\section{Non Redundant Parameterization of Pure States Adapted to ${\frak B}%
_{vec}$}

In order to produce non redundant parameterizations of pure states in ${\cal %
B}_{1}\left( {\cal H}^{N}\right) $ by the groups ${\frak G}_{fib}$ and $%
{\frak G}_{den}$ one needs to introduce biorthogonal bases for $K_{OD},K_{D}$
and $S_{D}$. (The space of trace class operators ${\cal B}_{1}\left( {\cal H}%
^{N}\right) $ is not a Hilbert space thus is not self dual, its dual space
is the space of bounded operators ${\cal B}\left( {\cal H}^{N}\right) )$.The
basis for $K_{OD},K_{D}$ and $S_{D}$ are 
\begin{eqnarray}
&&K_{OD}={\rm linear\,span}\left\{ \breve{h}_{\iota }:\breve{h}_{\iota
}=\int h_{\iota }\left( {\sf z}{\bf ,}{\sf z}^{\prime }\right) \left| {\sf z}%
\right\rangle \left\langle {\sf z}^{\prime }\right| d{\sf z}d{\sf z}^{\prime
};{\sf z}{\bf \neq }{\sf z}^{\prime }\right\} \\
&&K_{D}={\rm linear\,span}\left\{ \breve{k}_{\iota }:\breve{k}_{\iota }=\int
k_{\iota }\left( {\sf z}\right) \left| {\sf z}\right\rangle \left\langle 
{\sf z}\right| d{\sf z}\right.  \nonumber \\
&&\left. \int k_{\iota }\left( \kappa _{2,}..,\kappa _{N_{e}},\kappa
_{N_{e}+1},..\kappa _{N}\right) d\kappa _{2}\ldots d\kappa _{N_{e}}d\kappa
_{N_{e}+2}\ldots d\kappa _{N}{\bf =0}\text{ }{\bf a.e.}\right\}
\label{kbasis} \\
&&S_{D}={\rm linear\,span}\left\{ \breve{s}_{\iota }:\breve{s}_{\iota }=\int
s_{\iota }\left( {\sf z}\right) \left| {\sf z}\right\rangle \left\langle 
{\sf z}\right| d{\sf z}\right. ;  \nonumber \\
&&\left. \int s_{\iota }\left( \kappa _{2,}..,\kappa _{N_{e}},\kappa
_{N_{e}+1},..\kappa _{N}\right) d\kappa _{2}\ldots d\kappa _{N_{e}}d\kappa
_{N_{e}+2}\ldots d\kappa _{N}{\bf \neq 0\,a.e.}\right\}  \label{sbasis}
\end{eqnarray}%
where ${\bf a.e.\equiv }$ ''almost everywhere'' i.e. except on sets of
measure zero. These bases are biothorgonal to the bases $\left\{ \breve{k}%
_{\kappa }^{d},\breve{s}_{\kappa }^{d};1\leq \kappa \leq \infty \right\} $
of ${\cal B}\left( {\cal H}^{N}\right) $ i.e 
\begin{eqnarray}
\left( \breve{h}_{\kappa }^{d}|\breve{h}_{\iota }\right) &=&\left( \breve{k}%
_{\kappa }^{d}|\breve{k}_{\iota }\right) =\left( \breve{s}_{\kappa }^{d}|%
\breve{s}_{\iota }\right) =\delta _{\kappa \iota }  \nonumber \\
\left( \breve{k}_{\kappa }^{d}|\breve{s}_{\iota }\right) &=&\left( \breve{s}%
_{\kappa }^{d}|\breve{k}_{\iota }\right) =\left( \breve{h}_{\kappa }^{d}|%
\breve{s}_{\iota }\right)  \nonumber \\
&=&\left( \breve{h}_{\kappa }^{d}|\breve{k}_{\iota }\right) =\left( \breve{k}%
_{\kappa }^{d}|\breve{h}_{\iota }\right) =\left( \breve{s}_{\kappa }^{d}|%
\breve{h}_{\iota }\right) =0
\end{eqnarray}%
and we have introduced the super operator scalar product $\left( A|B\right)
=Tr\left\{ A^{\dagger }B\right\} $.

If we are given a reference pure state $D_{ref}$ that has a density $\gamma $
then it can be decomposed as 
\begin{equation}
D_{ref}^{N}{}_{K_{OD}}+D_{ref}^{N}{}_{K_{D}}+\Gamma \left( \gamma \right)
\end{equation}%
and we can define 
\begin{eqnarray}
\breve{h}_{1} &=&D_{ref}^{N}{}_{K_{OD}}  \nonumber \\
\breve{k}_{1} &=&D_{ref}^{N}{}_{K_{D}}  \nonumber \\
\breve{s}_{1} &=&\Gamma \left( \gamma \right)
\end{eqnarray}%
Then transformations that parameterize all other pure states in the fiber $%
\left[ \gamma \right] _{N}$ in terms of this reference and ${\frak G}_{fib}$
in a non redundant fashion can be produced from functions $\phi _{fib}$
defined as 
\begin{eqnarray}
\phi _{fib}\left( {\sf z}\right) &=&\sum_{1\leq \kappa \leq \infty }\phi
_{fib\kappa }\left( \left| {\sf z}\right\rangle \left\langle {\sf z}\right| |%
\breve{k}_{\kappa }\right) \left( D_{refK_{D}}^{Nd}|\left| {\sf z}%
\right\rangle \left\langle {\sf z}\right| \right)  \nonumber \\
\phi _{fib\kappa } &=&\int k_{\kappa }^{d\ast }\left( {\sf z}\right) \phi
_{fib}\left( {\sf z}\right) D_{ref}^{N}{}_{K_{D}}\left( {\sf z}\right) d{\sf %
z};\quad 1\leq \kappa \leq \infty  \label{refkd}
\end{eqnarray}%
where they also satisfy 
\begin{eqnarray}
\int s_{\iota }^{d\ast }\left( {\sf z}\right) \phi _{fib}\left( {\sf z}%
\right) s_{\kappa }\left( {\sf z}\right) d{\sf z} &=&0\quad \forall \iota
,\kappa  \nonumber \\
\int k_{\iota }^{d\ast }\left( {\sf z}\right) \phi _{fib}\left( {\sf z}%
\right) s_{\kappa }\left( {\sf z}\right) d{\sf z} &=&0\quad \forall \iota
,\kappa  \nonumber \\
\int s_{\iota }^{d\ast }\left( {\sf z}\right) \phi _{fib}\left( {\sf z}%
\right) k_{\kappa }\left( {\sf z}\right) d{\sf z} &=&0\quad \forall \iota
,\kappa  \nonumber \\
\int k_{\iota }^{d\ast }\left( {\sf z}\right) \phi _{fib}\left( {\sf z}%
\right) k_{1}\left( {\sf z}\right) d{\sf z} &\neq &0\quad \text{for some }%
\iota  \label{phicond}
\end{eqnarray}%
The conditions in eq. $\left( \ref{phicond}\right) $ are needed to ensure
that the transformation induced by $\phi _{fib}\left( {\sf z}\right) $
belong to the defining representation of ${\cal G}_{fib}$ on $K_{D}$. We can
the use the diagonal operators $T\left( \phi _{fib},\omega \right) =\int
\phi _{fib}\left( {\sf z}\right) e^{i\omega \left( {\sf z}\right) }\left| 
{\sf z}\right\rangle \left\langle {\sf z}\right| d{\sf z}$, to define
congruence transformations belonging to ${\frak G}_{fib}$ as 
\begin{eqnarray}
&&\hat{T}_{fib}\left( \phi _{fib},\omega \right) Y=T\left( \phi
_{fib},\omega \right) YT\left( \phi _{fib},\omega \right) ^{\dagger }=e^{%
{\cal \breve{V}}\left( \phi _{fib}\right) +i\breve{\omega}}Ye^{{\cal \breve{V%
}}\left( \phi _{fib}\right) -i\breve{\omega}}  \nonumber \\
&=&\left( 1+\breve{\phi}_{fib}\right) ^{\frac{1}{2}}e^{i\breve{\omega}%
}Y\left( 1+\breve{\phi}_{fib}\right) ^{\frac{1}{2}}e^{-i\breve{\omega}%
}=\left( 1+\hat{\Phi}_{fib}\right) \widehat{e^{i\breve{\omega}}}Y
\end{eqnarray}%
where the integral kernel of $\hat{\Phi}_{fibK_{D}}$ of the projection $\hat{%
P}_{K_{D}}\hat{\Phi}_{fib}\hat{P}_{K_{D}}$ of $\hat{\Phi}_{fibK_{D}}$ on $%
K_{D}$ is defined as $\Phi _{fibK_{D}}\left( {\sf z}_{1}{\sf ,z}_{2}\right)
=\phi _{fib}\left( {\sf z}_{1}\right) \delta \left( {\sf z}_{1}-{\sf z}%
_{2}\right) $. On diagonal operators $Y_{d}=Y_{K_{D}}+\Gamma \left( \gamma
\right) $ the action of these transformations are 
\begin{equation}
\hat{T}_{fib}\left( \phi _{fib},\omega \right) Y_{d}=\left( 1+\hat{\Phi}%
_{fib}\right) Y_{d}=Y_{d}+\hat{\Phi}_{fibK_{D}}Y_{K_{D}}=Y_{d}+\breve{\phi}%
_{fib}Y_{K_{D}}
\end{equation}%
as $\hat{P}_{S_{D}}\hat{\Phi}_{fib}\hat{P}_{S_{D}}=\hat{P}_{K_{D}}\hat{\Phi}%
_{fib}\hat{P}_{S_{D}}=0$. The full intra fiber transformation can thus be
expressed as 
\begin{equation}
\hat{T}_{fib}\left( \phi _{fib},\omega \right) =e^{\sum_{1\leq \kappa \leq
\infty }\left\{ \Upsilon _{K_{OD}\kappa }\left| \breve{h}_{\kappa }\right)
\left( D_{refK_{OD}}^{Nd}\right| +\Upsilon _{K_{D}\kappa }\left| \breve{k}%
_{\kappa }\right) \left( D_{refK_{D}}^{Nd}\right| \right\} }  \label{parker}
\end{equation}%
where the coefficients $\Upsilon _{K_{D}\kappa }$ are determined by $\phi
_{fib}$ and $\omega $ by the following relationships%
\begin{eqnarray}
{\cal V}\left( \phi _{fib}\right) \left( {\sf z}\right) &=&\left( 1+\phi
_{fib}\left( {\sf z}\right) \right) ^{\frac{1}{2}}  \nonumber \\
\Upsilon \left( {\sf z}\right) &=&{\cal V}\left( \phi _{fib}\right) \left( 
{\sf z}\right) +i\omega \left( {\sf z}\right)  \nonumber \\
\Upsilon \left( {\sf z}\right) &=&\sum_{1\leq \kappa \leq \infty }\Upsilon
_{K_{D}\kappa }\left( \left| {\sf z}\right\rangle \left\langle {\sf z}%
\right| |\breve{k}_{\kappa }\right) \left( D_{refK_{D}}^{Nd}|\left| {\sf z}%
\right\rangle \left\langle {\sf z}\right| \right)  \nonumber \\
\Upsilon _{K_{D}\kappa } &=&\int k_{\kappa }^{d\ast }\left( {\sf z}\right)
\Upsilon \left( {\sf z}\right) D_{ref}^{N}{}_{K_{D}}\left( {\sf z}\right) d%
{\sf z}
\end{eqnarray}
and the coefficients $\Upsilon _{K_{OD}\kappa }$ are determined by $\phi
_{fib}$ and $\omega $. When acting on $D_{ref}^{N}$ the transformation $\hat{%
T}_{fib}\left( \phi _{fib},\omega \right) $ produces all other pure states
in $\left[ \gamma \right] _{N}$ in a unique fashion and thus parameterizes
the pure states in $\left[ \gamma \right] _{N}$ in a non redundant manner in
terms of $\left\{ \phi _{fib},\omega \right\} $. Note that transformations
belonging to ${\frak G}_{fib}$ never change the norm of density operators as
they leave $\Gamma \left( \gamma \right) $ unchanged.

The transformations that parameterize $S_{D}$ in a non redundant fashion
from a reference pure state $D_{ref}^{N}$ can always be constructed from
selected elements of ${\frak G}_{den}$ acting on $Y_{ref}$. A complete
parameterization of $S_{D}$ can be achieved in terms of pure states, as
given an arbitrary density $\gamma $ there is {\em always at least one} pure
state that contracts to that density. If we define $\phi _{den}$ as 
\begin{eqnarray}
\phi _{den}\left( {\sf z}\right) &=&\sum_{1\leq \kappa \leq \infty }\phi
_{den\kappa }\left( \left| {\sf z}\right\rangle \left\langle {\sf z}\right| |%
\breve{s}_{\kappa }\right) \left( \Gamma \left( \gamma \right) ^{d}|\left| 
{\sf z}\right\rangle \left\langle {\sf z}\right| \right)  \nonumber \\
\phi _{den\kappa } &=&\int s_{\kappa }^{d\ast }\left( {\sf z}\right) \phi
_{den}\left( {\sf z}\right) \Gamma \left( \gamma \right) \left( {\sf z}%
\right) d{\sf z};\quad 1\leq \kappa \leq \infty
\end{eqnarray}%
where they also satisfy 
\begin{eqnarray}
\int s_{\iota }^{d\ast }\left( {\sf z}\right) \phi _{den}\left( {\sf z}%
\right) s_{1}\left( {\sf z}\right) d{\sf z} &\neq &0\quad \text{for some }%
\iota  \nonumber \\
\int k_{\iota }^{d\ast }\left( {\sf z}\right) \phi _{den}\left( {\sf z}%
\right) s_{\kappa }\left( {\sf z}\right) d{\sf z} &=&0\quad \forall \iota
,\kappa  \nonumber \\
\int s_{\iota }^{d\ast }\left( {\sf z}\right) \phi _{den}\left( {\sf z}%
\right) k_{\kappa }\left( {\sf z}\right) d{\sf z} &=&0\quad \forall \iota
,\kappa  \nonumber \\
\int k_{\iota }^{d\ast }\left( {\sf z}\right) \phi _{den}\left( {\sf z}%
\right) k_{1}\left( {\sf z}\right) d{\sf z} &=&0\quad \forall \iota ,\kappa
\end{eqnarray}%
We can the use the diagonal operators $T\left( \phi _{den}\right) =\int \phi
_{den}\left( {\sf z}\right) \left| {\sf z}\right\rangle \left\langle {\sf z}%
\right| d{\sf z}$, to define congruence transformations%
\begin{eqnarray}
\hat{T}_{den}\left( \phi _{den}\right) Y &=&T\left( \phi _{den}\right)
YT\left( \phi _{den}\right) ^{\dagger }=e^{{\cal \breve{V}}\left( \phi
_{den}\right) }Ye^{{\cal \breve{V}}\left( \phi _{den}\right) }  \nonumber \\
&=&\left( 1+\breve{\phi}_{den}\right) ^{\frac{1}{2}}Y\left( 1+\breve{\phi}%
_{den}\right) ^{\frac{1}{2}}=\left( 1+\hat{\Phi}_{den}\right) Y
\end{eqnarray}%
where the integral kernel of $\hat{\Phi}_{denS_{D}}$ of the projection $\hat{%
P}_{S_{D}}\hat{\Phi}_{den}\hat{P}_{S_{D}}$ of $\hat{\Phi}_{den}$ on $S_{D}$
is defined as $\Phi _{den}\left( {\sf z}_{1}{\sf ,z}_{2}\right) =\phi
_{den}\left( {\sf z}_{1}\right) \delta \left( {\sf z}_{1}-{\sf z}_{2}\right) 
$. On diagonal operators $Y_{d}=Y_{K_{D}}+\Gamma \left( \gamma \right) $
this transformation becomes 
\begin{equation}
\hat{T}_{den}\left( \phi _{den}\right) Y_{d}=\left( 1+\hat{\Phi}%
_{den}\right) Y_{d}=Y_{d}+\hat{\Phi}_{denS_{D}}\Gamma \left( \gamma \right)
=Y_{d}+\breve{\phi}_{den}\Gamma \left( \gamma \right)
\end{equation}%
as $\hat{P}_{K_{D}}\hat{\Phi}_{den}\hat{P}_{K_{D}}=\hat{P}_{K_{D}}\hat{\Phi}%
_{den}\hat{P}_{S_{D}}=0.$The full intra fiber transformation can thus be
expressed as 
\begin{equation}
\hat{T}_{den}\left( \phi _{den},\omega \right) =e^{\sum_{1\leq \kappa \leq
\infty }\left\{ \Upsilon _{K_{OD}\kappa }\left| \breve{h}_{\kappa }\right)
\left( \breve{Y}_{refK_{OD}}^{d}\right| +\Upsilon _{S_{D}\kappa }\left| 
\breve{s}_{\kappa }\right) \left( \Gamma \left( \gamma \right) ^{d}\right|
\right\} }  \label{refden}
\end{equation}%
where the coefficients $\Upsilon _{S_{D}\kappa }$ by the following
relationships%
\begin{eqnarray}
{\cal V}\left( \phi _{den}\right) \left( {\sf z}\right) &=&\left( 1+\phi
_{den}\left( {\sf z}\right) \right) ^{\frac{1}{2}}  \nonumber \\
\Upsilon \left( {\sf z}\right) &=&{\cal V}\left( \phi _{den}\right) \left( 
{\sf z}\right)  \nonumber \\
\Upsilon \left( {\sf z}\right) &=&\sum_{1\leq \kappa \leq \infty }\Upsilon
_{S_{D}\kappa }\left( \left| {\sf z}\right\rangle \left\langle {\sf z}%
\right| |\breve{s}_{\kappa }\right) \left( \Gamma \left( \gamma \right)
^{d}|\left| {\sf z}\right\rangle \left\langle {\sf z}\right| \right) 
\nonumber \\
\Upsilon _{K_{D}\kappa } &=&\int \breve{s}_{\kappa }^{d\ast }\left( {\sf z}%
\right) \Upsilon \left( {\sf z}\right) \Gamma \left( \gamma \right) \left( 
{\sf z}\right) d{\sf z}
\end{eqnarray}%
and the coefficients $\Upsilon _{K_{OD}\kappa }$ are determined by $\phi
_{den}$. When acting on $D_{ref}^{N}$ the transformation $\hat{T}%
_{den}\left( \phi _{den}\right) $ produces a unique set of pure states, one
in each equivalence class. Thus $\hat{T}_{den}\left( \phi _{den}\right) $
acting on the reference pure state $D_{ref}^{N}$ parameterizes $S_{D}$ in a
non redundant manner in terms of $\left\{ \phi _{den}\right\} $. It is
possible by left multiplying $\hat{T}_{den}\left( \phi _{den}\right) $ by a
transformation $\hat{T}_{fib}\left( \phi _{fib}\left( \phi _{den}\right)
,\omega \left( \phi _{den}\right) \right) $ uniquely determined by $\hat{T}%
_{den}\left( \phi _{den}\right) $, that a path of pure reference states $%
D_{ref}^{N}\left( \gamma \right) $ of a particular type (i.e. Independent
Particle, Antisymmetrized Geminal Power etc..) and obtain a more specific
non redundant parameterization of $S_{D}$.

It is further possible to choose a subset of the transformations given in
eq. $\left( \ref{refden}\right) $ that produce normalized pure states, (even
from an unnormalized $D_{ref}^{N}$), by letting $\Upsilon _{S_{D}1}$ depend
on $\left\{ \Upsilon _{S_{D}\kappa };2\leq \kappa \leq \infty \right\} $

\section{Production of Arbitrary Pure States from a Reference Pure State}

\noindent Here we prove that diagonal transformations, in the sense of
Appendix F, can produce any pure state from a reference pure state. The
equivalent mathematical statement is that such transformations are
transitive on the set of pure states.

Associated with any arbitrarily chosen pure state $|\Psi \rangle \,\in \,%
{\cal H}^{N}\,=\,{\cal H}^{N_{e}}\otimes {\cal H}^{N_{n}}$ there is a
wavefunction 
\begin{gather}
\Psi ({\bf x}_{1},\ldots ,{\bf x}_{N_{e}}{\bf X}_{1},\ldots ,{\bf X}%
_{N_{n}})=\langle {\bf x}_{1},\ldots ,{\bf x}_{N_{e}}{\bf X}_{1},\ldots ,%
{\bf X}_{N_{n}}|\Psi \rangle  \nonumber \\
=|\Psi ({\bf x}_{1},\ldots ,{\bf x}_{N_{e}}{\bf X}_{1},\ldots ,{\bf X}%
_{N_{n}})|e^{i\theta _{\Psi }({\bf x}_{1},\ldots ,{\bf x}_{N_{e}}{\bf X}%
_{1},\ldots ,{\bf X}_{N_{n}})}  \nonumber \\
\equiv e^{{\cal V}_{\Psi }({\bf x}_{1},\ldots ,{\bf x}_{N_{e}}{\bf X}%
_{1},\ldots ,{\bf X}_{N_{n}})}e^{i\theta _{\Psi }({\bf x}_{1},\ldots ,{\bf x}%
_{N_{e}}{\bf X}_{1},\ldots ,{\bf X}_{N_{n}})}  \label{psipolar}
\end{gather}

Consider $T_{d}\in {\cal C}({\cal H}^{N})$, a diagonal transformation in the
sense of Appendix F. Its integral kernel can be written as in eq. (\ref%
{Tdexplicit}) which leads to eq. (\ref{Tdcoord}). The latter equation in
combination with the last RHS of (\ref{psipolar}) gives 
\begin{equation}
\langle {\sf x},{\sf X}|T_{d}\Psi \rangle =e^{({\cal V}_{\Psi }+ {\cal V}%
)}e^{i(\theta_{\Psi }+\omega )}
\end{equation}%
with the obvious string of arguments in compressed notation or surpressed
for clarity.

If one chooses an arbitrary $|\chi \rangle \,\in \,{\cal H}^{N}$, then the
choice 
\begin{eqnarray}
{\cal V} &=&{\cal V}_{\chi }-{\cal V}_{\Psi }  \nonumber \\
\omega &=&\theta _{\chi }-\theta _{\Psi }  \label{VTthetaT}
\end{eqnarray}%
generates that $|\chi \rangle $ from the initial $|\Psi \rangle $, itself
arbitrarily chosen, via the transformation $T_{d}$. Since all physical
densities correspond to pure states \cite{Harriman81}, the result (\ref%
{VTthetaT}) means that all densities (fibers) can be reached given a pure
state corresponding to a reference density (fiber) via transformations of
the type developed in Appendix F.

\section{Partitioned Equation of Motion in the $\Upsilon $ Parameterization}

This Appendix tracks the reasoning given in Appendix A but for the complex
analytic parameterization, $\Upsilon ={\cal V}+i\omega $ of the fibers, eq. (%
\ref{Ffactor}). In this parameterization, the equations of motion that
follow from eq. (\ref{dEdlambda}) - (\ref{Eoft}) can be written in matrix
partitioned form as 
\begin{equation}
\left[ 
\begin{array}{ccc}
\eta _{{\cal \Upsilon }\,{\cal \Upsilon }}(t) & \eta _{{\cal \Upsilon
\Upsilon }^{\ast }}(t) & \eta _{{\cal \Upsilon }\gamma }(t) \\ 
\eta _{{\cal \Upsilon }^{\ast }{\cal \Upsilon }}(t) & \eta _{{\cal \Upsilon }%
^{\ast }{\cal \Upsilon }^{\ast }}(t) & \eta _{{\cal \Upsilon }^{\ast }\gamma
}(t) \\ 
\eta _{\gamma {\cal \Upsilon }}(t) & \eta _{\gamma {\cal \Upsilon }^{\ast
}}(t) & \eta _{\gamma \gamma }(t)%
\end{array}%
\right] \left[ 
\begin{array}{c}
{\cal \Upsilon } \\ 
{\cal \Upsilon }^{\ast } \\ 
\dot{\gamma}%
\end{array}%
\right] =\left[ 
\begin{array}{c}
\nabla _{{\cal \Upsilon }}{\cal E}\left( t\right) \\ 
\nabla _{{\cal \Upsilon }^{\ast }}{\cal E}\left( t\right) \\ 
\nabla _{\gamma }{\cal E}\left( t\right)%
\end{array}%
\right]  \label{3x3partmat}
\end{equation}%
where $^{\ast }$ denotes complex conjugation. These equations may be written
more compactly as 
\begin{equation}
\left[ 
\begin{array}{cc}
\eta _{{\bf \digamma \digamma }}\left( t\right) & \eta _{{\bf \digamma }%
\gamma }\left( t\right) \\ 
\eta _{\gamma {\bf \digamma }}\left( t\right) & \eta _{\gamma \gamma }\left(
t\right)%
\end{array}%
\right] \left[ 
\begin{array}{c}
{\bf \dot{\digamma}} \\ 
\dot{\gamma}%
\end{array}%
\right] =\left[ 
\begin{array}{c}
\nabla _{{\bf \digamma }}{\cal E}\left( t\right) \\ 
\nabla _{\gamma }{\cal E}\left( t\right)%
\end{array}%
\right]  \label{partmat2}
\end{equation}%
with 
\begin{eqnarray}
&&{\bf \digamma }=\left( 
\begin{array}{c}
{\cal \Upsilon } \\ 
{\cal \Upsilon }^{\ast }%
\end{array}%
\right) \,;\;\eta _{\digamma \digamma }\left( t\right) =\left( 
\begin{array}{cc}
\eta _{{\cal \Upsilon \Upsilon }} & \eta _{{\cal \Upsilon \Upsilon }^{\ast }}
\\ 
\eta _{{\cal \Upsilon }^{\ast }{\cal \Upsilon }} & \eta _{{\cal \Upsilon }%
^{\ast }{\cal \Upsilon }^{\ast }}%
\end{array}%
\right) \,;\;\eta _{{\bf \digamma }\gamma }\left( t\right) =\left( 
\begin{array}{c}
\eta _{{\cal \Upsilon }^{\ast }\gamma } \\ 
\eta _{{\cal \Upsilon }^{\ast }\gamma }%
\end{array}%
\right)  \nonumber \\
&&\eta _{\gamma {\bf \digamma }}\left( t\right) =\left( 
\begin{array}{cc}
\eta _{\gamma {\cal \Upsilon }} & \eta _{\gamma {\cal \Upsilon }^{\ast }}%
\end{array}%
\right) \,;\;\nabla _{\digamma }{\cal E}\left( t\right) =\left( 
\begin{array}{c}
\nabla _{{\cal \Upsilon }}{\cal E}\left( t\right) \\ 
\nabla _{{\cal \Upsilon }^{\ast }}{\cal E}\left( t\right)%
\end{array}%
\right)
\end{eqnarray}

In this ${\bf {\digamma }}$, $\gamma $ notation the formal structure of the
EOM eq. (\ref{partmat2}) is exactly as in eq. (\ref{partmat1}) of Appendix A
with corresponding definitions for $S$ and ${\cal E}$. In parallel with the
development leading to eq. (\ref{cdotgammadot}), it follows that 
\begin{eqnarray}
\dot{{\bf \digamma }}(t) &=&\eta _{{\bf \digamma \,\digamma }}^{-1}(t)\left[
\nabla _{{\bf \digamma }}{\cal E}(t)-\eta _{{\bf \digamma }\,\gamma }(t)\dot{%
\gamma}\right]  \nonumber \\
\dot{\gamma}(t) &=&\left[ \eta _{\gamma \,\gamma }(t)-\eta _{\gamma \,{\bf %
\digamma }}(t)\eta _{{\bf \digamma }\,{\bf \digamma }}^{-1}(t)\eta _{{\bf %
\digamma }\,\gamma }(t)\right] ^{-1}  \nonumber \\
&&\times \left[ \nabla _{\gamma }{\cal E}(t)-\eta _{\gamma \,{\bf \digamma }%
}(t)\eta _{{\bf \digamma }\,{\bf \digamma }}^{-1}(t)\nabla _{{\bf \digamma }}%
{\cal E}(t)\right]
\end{eqnarray}%
Under the assumption that $\eta _{{\bf \digamma \digamma }}^{-1}(t)$ exists,
denote the solution of the first of these for the correct density $\gamma
_{\#}(t)$ by ${\bf \digamma }_{\#}(t)$. Then the exact equation of motion
for the temporal development of the density is 
\begin{equation}
\dot{\gamma}(t)={\cal N}_{1}(\gamma (t))\left[ \nabla _{\gamma }{\cal E}%
(\gamma (t))-{\cal N}_{2}(\gamma (t))\right]
\end{equation}%
The functionals are defined as 
\begin{align}
{\cal N}_{1}(\gamma (t))& =\left[ \eta _{\gamma \,\gamma }({\bf \zeta }%
_{\#}(t)\,,\gamma (t))-\eta _{\gamma \,{\bf \digamma }}({\bf \zeta }%
_{\#}(t)\,,\gamma (t))\eta _{{\bf \digamma }\,{\bf \digamma }}^{-1}({\bf %
\zeta }_{\#}(t)\,,\gamma (t))\right. \times  \nonumber \\
& \qquad \qquad \qquad \qquad \qquad \qquad \qquad \qquad \qquad \qquad
\left. \eta _{{\bf \digamma }\,\gamma }({\bf \zeta }_{\#}(t)\,,\gamma (t))%
\right] ^{-1} \\
{\cal N}_{2}\left( \gamma (t)\right) & =\eta _{\gamma \,{\bf \digamma }}(%
{\bf \zeta }_{\#}(t)\,,\gamma (t))\eta _{{\bf \digamma }\,{\bf \digamma }%
}^{-1}({\bf \zeta }_{\#}(t)\,,\gamma (t))\nabla _{{\bf \digamma }}{\cal E}%
(\gamma (t)).\quad \\
{\cal E}(\gamma (t))& =\frac{\left\langle \Psi ({\bf \zeta }_{\#}(t),\gamma
(t))|H|\Psi ({\bf \zeta }_{\#}(t),\gamma (t))\right\rangle }{\left\langle
\Psi ({\bf \zeta }_{\#}(t),\gamma (t))|\Psi ({\bf \zeta }_{\#}(t),\gamma
(t))\right\rangle }
\end{align}

\begin{references}
\bibitem{Teil92} J.~Teilhaber, Phys. Rev. B {\bf 46}, 12990 (1992).

\bibitem{Stener95} M.~Stener, P.~Decleva, and A.~Lisini, J. Phys. B [At.
Mol. Opt. Phys.] {\bf 28}, 4973 (1995).

\bibitem{Saalmann96} U.~Saalmann and R.~Schmidt, Z. Phys. D {\bf 38}, 153
(1996).

\bibitem{Tong97} X.M.~Tong and S-I.~Chu, Phys. Rev. A {\bf 55}, 3406 (1997).

\bibitem{Serra97a} L.~Serra and A.~Rubio, Z. Physik D {\bf 40}, 262 (1997).

\bibitem{Ullrich97} C.A.~Ullrich, P.G.~Reinhard, and E.~Suraud, J. Phys.B
[At. Mol. Optical Phys.] {\bf 30}, 5043 (1997).

\bibitem{AG97} A.~G\"orling, Phys. Rev. B {\bf 55}, 2630 (1997).

\bibitem{Vignale97} G.~Vignale, C.A.~Ullrich, and S.~Conti, Phys. Rev. Lett. 
{\bf 79}, 4878 (1997).

\bibitem{Calvayrac98} F.~Calvayrac, A.~Domps, P.G.~Reinhard, E.~Suraud, and
C.A.~Ullrich, Euro. Phys. J. {\bf 4}, 207 (1998).

\bibitem{Reinhard98} P.G.~Reinhard, E.~Suraud, and C.A.~Ullrich, Euro. Phys.
J. {\bf 1}, 303 (1998).

\bibitem{Ullrich98} C.A.~Ullrich, P.G.~Reinhard, and E.~Suraud, Phys. Rev. A 
{\bf 57}, 1938 (1998).

\bibitem{Telnov98} D.A.~Telnov and S-I.~Chu, Phys. Rev. A {\bf 58}, 4749
(1998).

\bibitem{Tong98} X.-M.~Tong and S.-I.~Chu, Phys. Rev. A {\bf 57}, 452 (1998).

\bibitem{Campillo99} I.~Campillo, A.~Rubio, and J.M.~Pitarke, Phys. Rev.B 
{\bf 59}, 12188 (1999).

\bibitem{Vasiliev99} I.~Vasiliev, S.~Ogut, and J.R.~Chelikowsky, Phys. Rev.
Lett. {\bf 82}, 1919 (1999).

\bibitem{Knospe99} O.~Knospe, J.~Jellinek, U.~Saalmann, and R.~Schmidt, Eur.
Phys. J. D {\bf 5}, 1 (1999) and references therein.

\bibitem{Hessler} P.~Hessler, J.~Park, and K.~Burke, Phys. Rev. Lett. {\bf 82%
}, 378 (1999).

\bibitem{Gross96} E.K.U.~Gross, J.F.~Dobson, and M.~Petersilka, in {\it %
Density Functional Theory II - Topics in Current Chemistry 181}, R.F.
Nalewajski ed. (Springer, Berlin and Heidelberg, 1996), 81.

\bibitem{Casida95} M.~Casida in D.P. Chong ed., {\it Recent Advances in
Density Functional Methods} (World Scientific, Singapore, 1995), 155.

\bibitem{Bartolotti81} L.J.~Bartolotti, Phys. Rev. A {\bf 24}, 1661 (1981).

\bibitem{DebGhosh82} B.M.~Deb and S.K.~Ghosh, J. Chem. Phys. {\bf 77}, 342
(1982).

\bibitem{RG84} E.~Runge and E.K.U.~Gross, Phys. Rev. Lett. {\bf 52}, 997
(1984).

\bibitem{HK} P.~Hohenberg and W.~Kohn, Phys. Rev. {\bf 136}, B864 (1964).

\bibitem{Xu85} B-X.~Xu and A.K.~Rajagopal, Phys. Rev. A {\bf 31}, 2682
(1985).

\bibitem{Dhara87} A.K.~Dhara and S.K.~Ghosh, Phys. Rev. A {\bf 35}, 442
(1987).

\bibitem{Levy79} M.~Levy, Proc. Natl. Acad. Sci. USA {\bf 76}, 6062 (1979).

\bibitem{LevyPerdew85} M.~Levy and J.P.~Perdew in R.M.~Dreizler and J.~da
Providencia eds., {\it Density Functional Methods in Physics} (Plenum, N.Y.,
1985), 11.

\bibitem{Lieb} E.H.~Lieb, Internat. J. Quantum. Chem. {\bf 24}, 243 (1983).

\bibitem{KD86} H.~Kohl and R.M.~Dreizler, Phys. Rev. Lett. {\bf 56}, 1993
(1986).

\bibitem{Raj94A} A.K.~Rajagopal, Phys. Lett. A {\bf 193}, 363 (1994).

\bibitem{RajBuot94B} A.K.~Rajagopal and F.A.~Buot, Phys. Rev. E {\bf 50},
721 (1994).

\bibitem{LiTong86} T-C.~Li and P-q. Tong, Phys. Rev. A {\bf 34}, 529 (1986).

\bibitem{Qian98} A.~Qian and V.~Sahni, Phys. Lett. A {\bf 247}, 303 (1998).

\bibitem{Chermyak00} V.~Chernyak and S.~Mukamel J. Chem. Phys. {\bf 112},
3572 (2000).

\bibitem{DeumensA} E.~Deumens, A.~Diz, R.~Longo, and Y.~{\"O}hrn, Rev. Mod.
Phys. {\bf 66}, 917 (1994).

\bibitem{DeumensB} E.~Deumens and Y.~{\"O}hrn, J. Chem. Soc., Faraday Trans. 
{\bf 93}, 919 (1997) and refs.~therein.

\bibitem{BLWSBT99a} B.~Weiner and S.B.~Trickey, Adv. Quantum Chemistry {\bf %
35} (Academic, San Diego, 1999), 217.

\bibitem{BLWSBT00a} B.~Weiner and S.B.~Trickey, J. Molec. Struct. (Theochem) 
{\bf 501}-{\bf 02}, 65 (2000).

\bibitem{Leeuwen99} R.~van Leeuwen, Phys. Rev. Lett. {\bf 82}, 3863 (1999).

\bibitem{HEGexample} Even with a well-behaved temporal Taylor series, the
critical surface integral in the RG argument is a vulnerability. The jellium
model of a semi-infinite solid is an example. Place the jellium in the lower
half space with its surface at $z=0$, hence vacuum for $z>0$. Impose a
spatially and temporally varying electric field via some potential $v({\bf r}%
,t)$ that has a non-singular temporal Taylor series. Then pick another
potential $v^\prime$ such that the difference is 
\begin{eqnarray}
w({\bf r},t) & = & v({\bf r},t) - v^{\prime}({\bf r},t)  \nonumber \\
& \propto & E \, z \, cos(\omega \, t + \phi)  \nonumber
\end{eqnarray}
Clearly $v^\prime$ has a well-behaved Taylor series. The surface integral
that must be zero in the RG logic is 
\[
{\cal S} = \int d{\bf S}\cdot ( n_{0} \, w_{k} \, \nabla w_{k} ) %\nonumber
\]
with $w_{k}({\bf r})$ the {\it k}th partial time derivative of $w$ at $t=t_0$
. Clearly ${\cal S}$ is non-zero for the lower half-sphere for this example.
This kind of potential was disallowed by Dhara and Ghosh \cite{Dhara87} even
though it is important.

\bibitem{bareTaylor} A simple example is a single point charge $Q$ moving
along the $x$-axis with coordinate $X(t)$. Then 
\begin{eqnarray}
v({\bf r},t) & = & \frac{Q}{D^{1/2}(X(t),{\bf r})}  \nonumber \\
D(X(t),{\bf r}) & \equiv & (X(t) -x)^{2} + y^{2} + z^{2}  \nonumber
\end{eqnarray}
The leading singularity at second and third order is $D^{-3/2}(X(t_{0}),{\bf %
r})$, at fourth order $D^{-5/2}(X(t_{0}),{\bf r})$ and so on.

\bibitem{Xchneglect} Note also that a portion of the presentation in Ref. %
\citenum{Gross96} depends on the neglect of nuclear exchange other than
self-exchange. Therefore those specific developments are inapplicable to
atomic and molecular scattering.

\bibitem{KramSar} P.~Kramer and M.~Saraceno, {\it Geometry of the Time
Dependent Variational Principle in Quantum Mechanics} (Springer, New York
1981).

\bibitem{GrossKohn90} W.~Kohn and E.K.U.~Gross in {\it Density Functional
Theory for Many-Fermion Systems} S.B. Trickey spec. ed., Adv. Quant. Chem. 
{\bf 21} (Academic, San Diego, 1990), 255.

\bibitem{Fois88} E.S.~Fois, A.~Selloni, M.~Parinello,and R.~Car, J. Phys.
Chem. {\bf 92}, 3268 (1988).

\bibitem{PerdewSavin} J.P.~Perdew, A.~Savin, and K.~Burke, Phys. Rev. A {\bf %
51}, 4531 (1995).

\bibitem{BLWSBT98} B.~Weiner and S.B.~Trickey, Internat. J. Quantum Chem. 
{\bf 69}, 451 (1998).

\bibitem{SavinSem2} A.~Savin, in J.M. Seminario ed., {\it Recent
Developments and Applications of Modern Density Functional Theory}
(Elsevier, Amsterdam, 1996), 351.

\bibitem{spinpol} If necessary, an approximation corresponding to ordinary
spin-polarized DFT can be extracted from the fully spin-dependent theory by
neglecting the contributions of the off-diagonal elements of the
one-electron space-spin density kernel 
\begin{equation}
\left[ 
\begin{array}{cc}
D_{e,\alpha \alpha }^{1}\left( {\bf r,r}^{\prime }\right) & D_{e,\alpha
\beta }^{1}\left( {\bf r,r}^{\prime }\right) \\ 
D_{e,\beta \alpha }^{1}\left( {\bf r,r}^{\prime }\right) & D_{e,\beta \beta
}^{1}\left( {\bf r,r}^{\prime }\right)%
\end{array}%
\right] \approx \left[ 
\begin{array}{cc}
D_{e,\alpha \alpha }^{1}\left( {\bf r,r}^{\prime }\right) & 0 \\ 
0 & D_{e,\alpha \alpha }^{1}\left( {\bf r,r}^{\prime }\right)%
\end{array}%
\right]  \nonumber
\end{equation}

\bibitem{Sjoqvist} E.~Sj\"oqvist and M.~Hedstr\"om, Phys. Rev. A {\bf 56},
3417 (1997).

\bibitem{Bruhat} {\it Analysis, Manifolds, and Physics}, Y.~Choquet-Bruhat
and C. DeWitt-Morette with M. Dillard-Bleick, Revised Edition
(North-Holland, Amsterdam, 1982) 350 ff.

\bibitem{affinebundleref} {\it Encyclopedic Dictionary of Mathematics},
K.~It\^o ed. (MIT Press, Cambridge, 1996) {\bf 1}, 569.

\bibitem{Yref} For physical densities the existence of arbitrarily many
choices of $Y_{ref}$ is guaranteed by Harriman's theorem \cite{Harriman81}.
For the generalization to trace class operators, the guarantee follows from
the fact that any operator $Y$ always can be decomposed into four positive
operators $p_{j}$ summed as 
\[
Y=p_{1}-p_{2}+i(p_{3}-p_{4}) 
\]%
Under the contraction map $\Xi $, each of these yields a positive density
(not necessarily normed to the same value but that is no matter), hence each
of the $p_{j}$ can be constructed by application of Harriman's theorem.

\bibitem{Harriman81} J. Harriman, Phys. Rev. A {\bf 24}, 680 (1981).

\bibitem{diffbundle} Note that here we could consider a subtly different
fiber bundle with base the set of reference elements $\{Y_{ref}\}$ and fiber
elements 
\[
Y=\sum_{i=1}^{d}\alpha _{i}\xi _{i} 
\]%
We use this fiber bundle extensively in section V.

\bibitem{etainv} If $\eta^{-1}$ does not exist, classical equations of
motion nonetheless can be obtained in terms of a reduced number of variables
that produce a non-singular metric. For a discussion of proper treatment of
the zeros of $\eta(t,\, {\bf y},\, {\bf y}^{\prime}) $ in classical
mechanics, see, for example, {\it Classical Dynamics A Modern Perspective},
E.C.G.~Sudarshan and N.~Mukunda (John Wiley and Sons, NY, 1974) 112 ff.

\bibitem{BLW94} B.~Weiner, E.~Deumens, and Y.~\"{O}hrn, J. Math. Phys. {\bf %
35}, 1139 (1994).

\bibitem{congruence} A congruence transform of an operator A is of the form $%
TAT^{\dag }$

\bibitem{Ortiz81} J.V.~Ortiz, B.~Weiner, and Y.~\"Ohrn, Internat J. Quantum
Chem. S-{\bf 15}, 113 (1981).

\bibitem{Kummer} H.~Kummer, J. Math. Phys. {\bf 8}, 2063 (1967).

\bibitem{Blaizot86} J-P.~Blaizot and G.~Ripka, {\it Quantum Theory of Finite
Systems} (MIT Press, Cambridge MA, 1986) Chaps. 17 and 18 and references
therein.

\bibitem{BerezinShubin} F.A.~Berezin and M.A..~Shubin, {\it The
Schr\"odinger Equation} (Kluwer Academic Press, Dordrecht/Boston/London,
1981).

\bibitem{Richtmyer} R.~Richtmyer, {\it Principles of Advanced Mathematical
Physics} (Springer Verlag, NY, 1978) {\bf 1}, 245.

\bibitem{KreibGross} T. Kreibich and E.K.U. Gross, Phys. Rev. Lett {\bf 86},
2984 (2001).
\end{references}

\end{document}
